%% ==========================================================================
%% DRAFT ONLY -- NOT FOR SUBMISSION, POSTING, OR CIRCULATION.
%% Coordinate with W. G. P. Mayner before any public posting (see Section 1.4).
%% Author byline: Daniel Kirtchakov (Independent researcher).
%% Bibliographic details of arXiv:1802.01225, arXiv:2410.06959, and the
%%   Tao and Speyer blog posts verified against the sources on 2026-08-04.
%%
%% RE-SCOPED 2026-08-04 after a first-hand prior-art check. Sources read in
%% full and verified on 2026-08-04:
%%   * github.com/wmayner/dixmier-counterexample (REPORT.md, dixmier-note.tex;
%%     both commits 2026-07-21). REPORT.md Sec.8 contains the cotangent lift
%%     Phi(q,p)=(F(q),G(q)p), "preserves the standard symplectic form exactly",
%%     det DPhi = 1, non-injective; verified-facts table row 24; priority
%%     Sec.12 item 3 lists "The symplectic lift to C^6 (Sec.8)".
%%   * arXiv:2607.20210 (Shaska), source of v1 (2026-07-22) and v2
%%     (2026-07-25) downloaded and read. v2 Thm 8.3 = the master equation;
%%     v2 Lem 8.5 = the base-point/anchor identity; v1 Thm 6.1 = v2 Thm 7.1
%%     = order-two vanishing on the contracted locus; v1/v2 Thm 4.4(2) = S3
%%     monodromy; v1/v2 Prop 5.1 = fiber counts 3/1/0 and the missed curve
%%     {(4/27 t^-2, 4/3 t^-1, t)}; v2 Thm 10.10 = the (1,1,d3) emptiness only.
%%   * ulam.ai/research/jacobian.pdf, PDF CreationDate 2026-07-20 05:54 EDT.
%%     Thm 4.2 = fiber counts 3/1/0 and image complement Gamma; Thm 5.1 =
%%     the (lambda,a,c,H) family; Cor 5.3 = every generic degree >= 3.
%%   * aaronlou.com/jacobian_counterexample_derivation.pdf, 2026-07-20 04:32.
%%   * MathOverflow 513387 (2026-07-20 07:56 UTC, 0 answers) and 513392
%%     (2026-07-20 11:40 UTC, 0 answers). MO 513390 could not be retrieved
%%     (404 / absent from the API) and is NOT cited.
%% ==========================================================================
%% CHANGELOG
%%   v0.1.2  2026-08-06  ERRATUM (fiber count over {Delta_2 = 0}).
%%     Corrected Theorem D (statement, and the "Fiber counts" paragraph of its
%%     proof): the fiber of F over the pullback of {Delta_2 = 0} (with
%%     Delta_1 != 0) has THREE distinct points, not two. {Delta_2 = 0} is an
%%     apparent branch locus -- the invariant u = 1+xy fails to separate two of
%%     the three unramified sheets, which is why Delta_2 occurs SQUARED in
%%     disc = -4 Delta_1 Delta_2^2. The achievable set-theoretic fiber sizes
%%     are exactly {3, 1, 0}; the value 2 never occurs. Independent
%%     re-verification (exact arithmetic, written from the raw map): the
%%     scripts in scripts/erratum-check/ (fibre_check.py, exhibit2.py,
%%     structural.py). No other claim in the note depended on the erroneous
%%     value: the generic degree 3, the image im F = C^3 \ Gamma, and the S_3
%%     monodromy from disc = -4 Delta_1 Delta_2^2 are unaffected. Supersedes
%%     v0.1.0 and v0.1.1 on this point only. See ERRATUM-v0.1.2.md.
%% ==========================================================================
\documentclass[11pt]{amsart}

\usepackage[margin=1.2in]{geometry}
\usepackage{amsmath,amssymb}
\usepackage{booktabs}
\usepackage[colorlinks=true,allcolors=blue]{hyperref}

\newcommand{\CC}{\mathbb{C}}
\newcommand{\QQ}{\mathbb{Q}}
\newcommand{\Cstar}{\CC^{\times}}
\newcommand{\W}{W}
\newcommand{\im}{\operatorname{im}}
\newcommand{\ord}{\operatorname{ord}}
\newcommand{\gr}{\operatorname{gr}}
\newcommand{\id}{\operatorname{id}}
\newcommand{\disc}{\operatorname{disc}}
\newcommand{\JC}{\mathrm{JC}}
\newcommand{\DC}{\mathrm{DC}}
\newcommand{\PC}{\mathrm{PC}}
\newcommand{\Jb}{\mathcal{J}}
\newcommand{\brk}{\mathcal{B}}

\theoremstyle{plain}
\newtheorem{theorem}{Theorem}[section]
\newtheorem{lemma}[theorem]{Lemma}
\newtheorem{proposition}[theorem]{Proposition}
\newtheorem{corollary}[theorem]{Corollary}
\newtheorem*{thmA}{Theorem A}
\newtheorem*{thmB}{Theorem B}
\newtheorem*{thmC}{Theorem C}
\newtheorem*{thmD}{Theorem D}
\newtheorem*{thmE}{Theorem E}
\theoremstyle{definition}
\newtheorem{definition}[theorem]{Definition}
\theoremstyle{remark}
\newtheorem{remark}[theorem]{Remark}

\begin{document}

\title[Degree minimality and the moment map for the Alp\"oge Keller
map]{Degree minimality in the equivariant class of the Alp\"oge Keller map,
and the moment-map structure of its cotangent lift}

\author{Daniel Kirtchakov}
\address{Independent researcher}
\email{daniel@halfounce.io}
\thanks{All computations in this note were carried out with Claude (Fable~5)
driving SymPy~1.14 and msolve~0.10.1. Every displayed identity is asserted by
a script listed in Section~\ref{sec:repro}; nothing is conjectural unless
labeled so. \emph{Draft of August 4, 2026; erratum release v0.1.2 of
August 6, 2026 (Theorem D fiber count over $\{\Delta_{2}=0\}$; see the
footnote there) --- not for circulation.}}

\subjclass[2020]{Primary 14R15; Secondary 16S32, 17B63, 53D55}
\keywords{Jacobian conjecture, Dixmier conjecture, Poisson conjecture,
Weyl algebra, Keller map, polynomial symplectomorphism}

\date{August 4, 2026}

\begin{abstract}
Let $F\colon\CC^{3}\to\CC^{3}$ be the degree-$7$ Keller map announced by
L.~Alp\"oge on July 19, 2026, which refutes the Jacobian conjecture in
dimension three. Our two new results concern its $\Cstar$-equivariant
class, in which $F=(A/x^{2},B/x,xC)$ with $A,B,C$ polynomials in the
invariants $u=1+xy$, $s=x^{2}z$.

First, a \emph{degree-minimality theorem}: in that class no Keller map of
degree $\le6$ exists at all in the sector where $C$ genuinely involves $s$
--- the sector containing Alp\"oge's map --- and every Keller map of
degree $\le6$ elsewhere in the class is a polynomial automorphism of
$\CC^{3}$. Degree $7$ is therefore minimal in the class. The emptiness
statements rest on eight Gr\"obner-basis unit-ideal certificates computed
over $\QQ$ (msolve, multi-modular F4) and reproduced mod $32003$, together
with a \emph{no-go lemma} valid for every weight $k\ge1$ and every degree:
if $C$ is $s$-free and $A,B$ are at most affine in $s$, the map is an
automorphism.

Second, a \emph{moment-map identity}. The $\Cstar$-action lifts
Hamiltonianly to $T^{*}\CC^{3}$, and the cotangent lift
$\Phi(q,p)=\bigl(F(q),(JF(q))^{-\mathsf T}p\bigr)$ preserves the moment
map exactly: $\nu\circ\Phi=\mu$ with $\mu=xp_{1}-yp_{2}-2zp_{3}$.

The map $\Phi$ itself is not ours: it was written down, and its exact
preservation of the symplectic form, $\det J\Phi=1$, and non-injectivity
verified, by W.~G.~P.~Mayner on July 21, 2026. We identify it as an
explicit witness for the Poisson conjecture $\PC_{3}$ via the
Adjamagbo--van~den~Essen chain, record its component degrees
$(7,6,4,9,10,12)$ and an explicit rational triple collision in $\CC^{6}$,
and prove the quantization identity $\gr\Psi_{F}=\Phi^{*}$ for Mayner's
Weyl-algebra endomorphism $\Psi_{F}$. The remaining structural material
--- the master equation, the base-point (anchor) identity, the $S_{3}$
cover, and the exact image theorem --- was obtained independently here but
is anticipated by Shaska, by the anonymous ulam.ai note, and by Mayner,
all in July 2026; it is retained for self-containedness and priority is
credited to those sources in Section~\ref{ssec:credit}. Everything is in
characteristic zero; $\JC_{2}$, $\DC_{1}$, and $\DC_{2}$ remain open.
\end{abstract}

\maketitle

\section{Introduction}\label{sec:intro}

\subsection{The three conjectures}
Throughout, we work over $\CC$; every statement in this note is in
characteristic zero. For a polynomial map $F=(F_{1},\dots,F_{n})\colon
\CC^{n}\to\CC^{n}$ write $JF=(\partial F_{i}/\partial x_{j})$ for its
Jacobian matrix. $F$ is a \emph{Keller map} if $\det JF$ is a nonzero
constant \cite{Keller39}. Three long-standing conjectures assert that
natural classes of endomorphisms are automorphisms:
\begin{itemize}
\item $\JC_{n}$ (Jacobian conjecture): every Keller map
$\CC^{n}\to\CC^{n}$ is a polynomial automorphism.
\item $\DC_{n}$ (Dixmier conjecture \cite{Dixmier68}): every unital
$\CC$-algebra endomorphism of the $n$-th Weyl algebra $\W_{n}$ is an
automorphism.
\item $\PC_{n}$ (Poisson conjecture): every unital $\CC$-algebra
endomorphism of the polynomial Poisson algebra
$P_{n}=\bigl(\CC[q_{1},\dots,q_{n},p_{1},\dots,p_{n}],\{\,,\}\bigr)$
that preserves the Poisson bracket is an automorphism.
\end{itemize}
These are linked by the chain
\begin{equation}\label{eq:chain}
\JC_{2n}\;\Longrightarrow\;\PC_{n}\;\Longrightarrow\;\DC_{n}
\;\Longrightarrow\;\JC_{n},
\end{equation}
established by Tsuchimoto \cite{Tsuchimoto05}, Belov-Kanel--Kontsevich
\cite{BKK07}, Bavula \cite{Bavula05}, and Adjamagbo--van den Essen
\cite{AvdE07} (who inserted the Poisson link); for the implication
$\DC_{n}\Rightarrow\JC_{n}$ see \cite{vdE00} and \cite[p.~2]{BKK07}.

On July 19, 2026, Levent Alp\"oge announced an explicit counterexample to
$\JC_{3}$ \cite{Alpoge26} --- the map $F$ of \eqref{eq:themap} below ---
thereby refuting $\JC_{n}$ for all $n\ge3$. (The problem had been suggested to him by Akhil
Mathew, and, per the announcement, the search that produced the map was run
with Claude Fable~5. We do not describe the discovery process beyond the
announcement; see Tao's expository post \cite{Tao26} and the discussion at
the Secret Blogging Seminar \cite{Speyer26}.) By the chain
\eqref{eq:chain}, read contrapositively, the failure of $\JC_{3}$ formally
entails the failure of $\DC_{3}$ and hence of $\PC_{3}$. Both consequences
are already on the public record: Mayner \cite{Mayner26} gave the explicit
Weyl-algebra witness for $\neg\DC_{3}$ on July 20--21, 2026, and the
Wikipedia article on the Jacobian conjecture states flatly that ``the
Dixmier and Poisson conjectures are false in every dimension $n>2$'',
citing \cite{AvdE07}. Nothing in Sections
\ref{sec:weyl}--\ref{sec:poisson} should be read as a priority claim on
either statement; what those sections contribute is an explicit
$\PC_{3}$-witness, its degrees and momenta, and the quantization identity
relating it to $\Psi_{F}$. The new results of this note are the
degree-minimality theorem (Theorem~E) and the moment-map identity
(Theorem~C(i)); see Section~\ref{ssec:credit} for a claim-by-claim
priority table.

\subsection{The map}
Alp\"oge's map is $F=(F_{1},F_{2},F_{3})\colon\CC^{3}\to\CC^{3}$,
\begin{equation}\label{eq:themap}
\begin{aligned}
F_{1}&=(1+xy)^{3}z+y^{2}(1+xy)(4+3xy), &&\deg F_{1}=7,\\
F_{2}&=y+3x(1+xy)^{2}z+3xy^{2}(4+3xy), &&\deg F_{2}=6,\\
F_{3}&=2x-3x^{2}y-x^{3}z, &&\deg F_{3}=4,
\end{aligned}
\end{equation}
with $\det JF\equiv-2$. It is non-injective in a way anyone can check by
hand in a minute:
\begin{equation}\label{eq:collision}
F\Bigl(1,-\tfrac32,\tfrac{13}2\Bigr)
=F\Bigl(-1,\tfrac32,\tfrac{13}2\Bigr)
=F\Bigl(0,0,-\tfrac14\Bigr)
=\Bigl(-\tfrac14,0,0\Bigr),
\end{equation}
a triple collision at rational points.

\subsection{Statement of results}\label{ssec:results}

\emph{What is new here is Theorem~E and Theorem~C(i)}, together with the
no-go Lemma~\ref{lem:nogo} on which Theorem~E rests. Theorem~E says that
in the $\Cstar$-equivariant class of Alp\"oge's map, degree $7$ is
minimal: below it the relevant sector is \emph{empty}, and what survives
elsewhere in the class is automorphisms only. Theorem~C(i) says that the
$\Cstar$-action lifts Hamiltonianly to the cotangent bundle and that the
cotangent lift preserves the moment map on the nose. The remaining
theorems are stated because the note is meant to be self-contained; each
carries its priority attribution in its own statement, and
Section~\ref{ssec:credit} gives the full table. The reader interested only
in the new material may go directly to Sections~\ref{ssec:moment} and
\ref{sec:min}.

Let $N=(JF)^{-1}=\operatorname{adj}(JF)/(-2)$, a $3\times3$ matrix with
\emph{polynomial} entries (of total degree at most $11$, at most $14$
monomials each --- $77$ in total across the nine entries), and let
$D_{j}=\sum_{k=1}^{3}N_{kj}\,\partial_{k}$ for $j=1,2,3$.

\begin{thmA}[Dixmier conjecture fails for $n\ge3$; \emph{Mayner's theorem},
independently verified here]
The assignments $x_{i}\mapsto F_{i}$, $\partial_{j}\mapsto D_{j}$ extend to
a well-defined unital $\CC$-algebra endomorphism
$\Psi_{F}\colon\W_{3}\to\W_{3}$. It is injective, and it is not surjective:
the generator $x$ is not in its image. Consequently $\DC_{n}$ is false for
every $n\ge3$, with explicit witness $\Psi_{F}\otimes\id_{\W_{n-3}}$.
\emph{This is due to W.~G.~P.~Mayner} (comment of July 20, 2026 on
\cite{Speyer26}; note and report of July 21, 2026 \cite{Mayner26}, where
the non-member exhibited is $y$ rather than $x$). We claim no priority
whatsoever for it; Section~\ref{sec:weyl} is an independent
re-verification, retained because Theorem~B and Theorem~C(ii) are stated
relative to $\Psi_{F}$.
\end{thmA}

\begin{thmB}[Mayner's cotangent lift, identified as a counterexample to
the Poisson conjecture]
\emph{The map $\Phi$ below and its three displayed properties are
Mayner's} \cite[\S8]{Mayner26}, who verified there that the cotangent lift
$\Phi(q,p)=(F(q),G(q)p)$, $G=(JF)^{-\mathsf T}$, is a polynomial self-map
of $\CC^{6}$ which ``preserves the standard symplectic form exactly'', has
$\det J\Phi=1$ identically, and is not injective; his priority section
lists ``the symplectic lift to $\CC^{6}$'' among the items apparently not
previously public. \emph{What is added here} is: the identification of
$\Phi^{*}$ as an explicit witness for $\PC_{n}$ through the
Adjamagbo--van~den~Essen formulation \cite{AvdE07}; the component degrees;
the explicit momenta of a rational triple collision; and, in Theorem~C,
the moment map and the quantization identity. Precisely:
\[
\Phi\colon\CC^{6}\to\CC^{6},\qquad
\Phi(q,p)=\bigl(F(q),\,(JF(q))^{-\mathsf T}p\bigr)
\]
is a polynomial map with component degrees $(7,6,4,9,10,12)$ whose
Jacobian $M=J\Phi$ satisfies $M^{\mathsf T}\Omega M=\Omega$ identically and
$\det M=1$, where $\Omega$ is the standard symplectic matrix. $\Phi$ is
generically $3{:}1$; in particular the three rational points
\[
a_{1}=\Bigl(1,-\tfrac32,\tfrac{13}2;\,-\tfrac{405}{16},\tfrac{25}{8},
-\tfrac{13}{8}\Bigr),\quad
a_{2}=\Bigl(-1,\tfrac32,\tfrac{13}2;\,-\tfrac{387}{16},\tfrac{31}{8},
\tfrac{11}{8}\Bigr),\quad
a_{3}=\Bigl(0,0,-\tfrac14;\,\tfrac92,2,1\Bigr)
\]
satisfy $\Phi(a_{1})=\Phi(a_{2})=\Phi(a_{3})=
\bigl(-\tfrac14,0,0;\,1,2,3\bigr)$. Consequently the pullback $\Phi^{*}$ is
an injective, non-surjective Poisson endomorphism of $P_{3}$, and $\PC_{n}$
is false for every $n\ge3$.
\end{thmB}

\begin{thmC}[Equivariant structure; \emph{part (i) is new}]
\leavevmode
\begin{enumerate}
\item[(i)] The $\Cstar$-action $t\cdot(x,y,z)=(tx,t^{-1}y,t^{-2}z)$ lifts
to a Hamiltonian action on $T^{*}\CC^{3}\cong\CC^{6}$ with moment map
$\mu=xp_{1}-yp_{2}-2zp_{3}$; the target carries the action of weights
$(-2,-1,1)$ with moment map $\nu=-2Q_{1}P_{1}-Q_{2}P_{2}+Q_{3}P_{3}$. The
map $\Phi$ is equivariant and preserves the moment map exactly:
$\nu\circ\Phi=\mu$.
\item[(ii)] $\Psi_{F}$ preserves the order filtration of $\W_{3}$, and
under the canonical isomorphism $\gr\W_{3}\cong\mathcal{O}(T^{*}\CC^{3})$
one has $\gr\Psi_{F}=\Phi^{*}$. In this precise sense $\Psi_{F}$ is the
canonical quantization of $\Phi$, and quantization does not restore
invertibility. (The symbol computation behind this --- that the
$\sigma_{1}(D_{j})=\sum_{k}N_{kj}\xi_{k}$ form a $\CC[x,y,z]$-basis of the
degree-one part of $\gr\W_{3}$ --- is step~2 of Mayner's self-contained
proof \cite[\S8]{Mayner26}; the identification of the resulting graded
endomorphism with $\Phi^{*}$ is our packaging.)
\end{enumerate}
The $\Cstar$-equivariance of $F$ itself is not new: it is recorded by
Speyer \cite{Speyer26}, in \cite{MO513387}, by Mayner
\cite[\S6]{Mayner26}, and is the organizing hypothesis of Shaska
\cite{Shaska26}. Part~(i) --- the Hamiltonian lift and the exact identity
$\nu\circ\Phi=\mu$ --- is, to our knowledge, not in any of them; we found
no occurrence of ``moment map'' anywhere in the July--August 2026
literature on this example.
\end{thmC}

\begin{thmD}[Exact image; \emph{anticipated}, see the attribution below]
The image of $F$ is $\CC^{3}$ minus the punctured rational curve
\[
\Gamma=\Bigl\{\Bigl(\tfrac{4}{27t^{2}},\,\tfrac{4}{3t},\,t\Bigr):
t\in\Cstar\Bigr\},
\]
and consequently $\im\Phi=\CC^{6}\setminus(\Gamma\times\CC^{3})$. Generic
fibers of $F$ have three points; over the pullback of $\{\Delta_{2}=0\}$
(with $\Delta_{1}\neq0$) the fiber still has three \emph{distinct}
points,\footnote{\emph{Erratum, v0.1.2 (2026-08-06).} Releases v0.1.0 and
v0.1.1 stated here that the count ``drops (generically) to two over the
pullback of $\{\Delta_{2}=0\}$,'' and the proof below read ``two sheets
merge.'' That was incorrect. The fiber over $\{\Delta_{2}=0\}$ (with
$\Delta_{1}\neq0$) has three distinct points; $\{\Delta_{2}=0\}$ is an
apparent branch locus. The achievable set-theoretic fiber sizes are exactly
$\{3,1,0\}$; the value $2$ never occurs. This was re-verified independently
in exact arithmetic from the raw map (scripts
\texttt{fibre\_check.py}, \texttt{exhibit2.py}, \texttt{structural.py} in
\texttt{scripts/erratum-check/}; see \texttt{ERRATUM-v0.1.2.md}). No other
statement in this note used the erroneous value: the generic degree $3$, the
image $\im F=\CC^{3}\setminus\Gamma$, and the $S_{3}$ monodromy of
Proposition~\ref{prop:cover}(ii) are unaffected.} since
$\{\Delta_{2}=0\}$ is an \emph{apparent} branch locus --- the invariant
$u=1+xy$ fails to separate two of the three unramified sheets, which is why
$\Delta_{2}$ occurs squared in $\disc=-4\Delta_{1}\Delta_{2}^{2}$
(Proposition~\ref{prop:cover}(ii)). The count drops to one over the pullback
of $\{\Delta_{1}=0\}$ --- two sheets escaping to infinity --- and to zero
exactly on $\Gamma$, where all three sheets have escaped. The achievable
set-theoretic fiber sizes are exactly $\{3,1,0\}$; the value $2$ does not
occur. ($\Delta_{1},\Delta_{2}$ are the explicit discriminant polynomials of
Proposition~\ref{prop:cover}.)

\emph{Priority.} The fiber counts $3/1/0$ and the identification of the
missed locus with the curve $\Gamma$ were obtained first in the anonymous
note \cite[Thm.~4.2]{Ulam26} (PDF timestamp July 20, 2026), and
independently in \cite[\S4.3]{Mayner26} and in \cite[Prop.~5.1]{Shaska26},
where the same curve appears in the same normalization,
$\{(\tfrac{4}{27}t^{-2},\tfrac43 t^{-1},t)\}$. Our proof is independent
and is retained only for self-containedness and because the transfer to
$\im\Phi$ is used later.
\end{thmD}

\begin{thmE}[Degree minimality in the equivariant class; \emph{new}]
In the equivariant class of weights $(1,-1,-2)$ --- maps
$F=(A/x^{2},\,B/x,\,xC)$ with $A,B,C\in\CC[u,s]$, $u=1+xy$, $s=x^{2}z$,
subject to the polynomial-lift conditions --- every Keller map of degree
$\le6$ is a polynomial automorphism of $\CC^{3}$. Stronger: in the entire
sector where $C$ genuinely involves $s$ (the sector containing Alp\"oge's
map), no Keller map of degree $\le6$ exists at all, automorphisms
included. Hence degree $7$ is minimal in this class. Unconditional
minimality of degree $7$ among all counterexamples in $\CC^{3}$ remains
open.
\end{thmE}

Theorem~E is proved in Section~\ref{sec:min}; it rests on the no-go
Lemma~\ref{lem:nogo} (valid for all weights $k\ge1$ and all degrees, and
also new) together with eight Gr\"obner-basis unit-ideal certificates
computed over $\QQ$ (msolve, multi-modular F4) and reproduced mod
$32003$. We record how it sits against the nearest published statements.

\begin{itemize}
\item Shaska \cite[Thm.~10.10]{Shaska26} proves that no graded Keller
counterexample in dimension three has \emph{both} $A$ and $B$ of degree
one in the invariants, for any signature $(r,s)$ --- a strictly weaker
emptiness statement, and one he explicitly leaves incomplete: ``the
smallest cases left open in dimension three are $(r,s)=(1,1)$ with
$\mathbf d=(2,1,d_{3})$ and $\mathbf d=(1,2,d_{3})$; whether
$\mathcal K(3,(1,-1,-1))$ is empty we do not know.'' Theorem~E clears an
entire degree band, in his hyperbolic signature $(r,s)=(1,2)$, with no
degree-one hypothesis on $A$ or $B$.
\item Mayner \cite[\S7]{Mayner26} proves a \emph{quadratic no-go} inside a
different ansatz --- the $z$-linear line congruences with a quadratic
direction field --- concluding that ``within the parallel-congruence
construction, degree $3$ is the minimum degree of the covering''. That is
a statement about the covering degree in his structural class; Theorem~E
is a statement about the total degree in the $\Cstar$-equivariant class.
Neither contains the other.
\item Jelonek \cite{Jelonek26} shows that if $X(n,d)$ is irreducible then
for $n\ge3$, $d\ge6$ a \emph{generic} element is a counterexample. That
is the opposite direction: an abundance statement above degree $6$, not an
emptiness statement below it, and it is conditional on irreducibility.
\item Unconditionally, invertibility of Keller maps is known only through
degree $2$, in every dimension, by Wang's theorem \cite{Wang80}; and
ulam.ai \cite[Cor.~5.3]{Ulam26} exhibits nonproper Keller maps of every
generic degree $d\ge3$ (their total degrees grow, and none of their
degree-$\le6$ members is equivariant of our type). Theorem~E is therefore
the first emptiness result we know of that rules out an explicit degree
band inside a class known to contain a counterexample.
\end{itemize}

The remaining proofs occupy Sections
\ref{sec:weyl}--\ref{sec:equiv}. Every displayed identity is asserted by a
script listed in Section~\ref{sec:repro}.

\subsection{Provenance and credit}\label{ssec:credit}
This subsection was rewritten on August 4, 2026 after a first-hand check of
the prior art, in which several results stated here turned out to have
earlier public sources. What follows is the corrected record.

\emph{The map.} The counterexample $F$ is Alp\"oge's, announced July 19,
2026; the problem was posed by Akhil Mathew, and the search was run with
Claude Fable~5, per the announcement. Everything in this note is downstream
of that example.

\emph{The Weyl endomorphism \underline{and} the symplectic lift.} Both are
W.~G.~P.~Mayner's, from the same day and the same source. The explicit
operators $D_{j}$ and the statement that $\Psi_{F}$ refutes $\DC_{3}$ were
made public in a comment of July 20, 2026 on the Secret Blogging Seminar
thread \cite{Speyer26} and in a six-page informal note of July 21, 2026
(prepared with Claude Fable~5); the accompanying report, in the same
repository and the same two commits of July 21, 2026, contains in
\S8 the cotangent lift $\Phi(q,p)=(F(q),G(q)p)$ together with the
statement --- machine-verified there --- that it preserves the standard
symplectic form exactly, that $\det J\Phi=1$, and that it is not
injective; his verified-facts table lists it as item~24, and his priority
section \S12 lists ``the symplectic lift to $\CC^{6}$'' as item~3 among
the results apparently not previously public \cite{Mayner26}. Mayner also
records the correct caution that this does \emph{not} bear on the
Belov-Kanel--Kontsevich automorphism conjecture, a caution we repeat in
Section~\ref{ssec:open}. An earlier draft of the present note credited
Mayner for the Weyl endomorphism only; that was an error of attribution,
and it is corrected here. Mayner's note additionally goes further than our
Section~\ref{sec:weyl}: it computes the exact image
$\bigoplus_{\alpha}\CC[F]D^{\alpha}$, exhibits non-members, tabulates the
codimension of the image in filtration degrees $d\le7$, and observes that
the cokernel is not finitely generated. We claim no priority for either
$\Psi_{F}$ or $\Phi$.

\emph{Independently obtained, subsequently found to be anticipated.} The
material of Sections~\ref{ssec:reduction}, \ref{ssec:cover},
\ref{ssec:image} and Lemma~\ref{lem:anchor} was obtained here
independently, before we became aware of the sources below; it is
\emph{not} new, and we retain it only for self-containedness. Priority
belongs to those sources, as follows. All dates are 2026 and were checked
against the sources themselves (arXiv submission records, PDF timestamps,
Git commit timestamps, StackExchange creation timestamps) on August 4,
2026.

\begin{center}
\small
\begin{tabular}{p{0.30\textwidth}p{0.42\textwidth}l}
\toprule
Item (as stated here) & Earliest public source we could verify & Date\\
\midrule
Master equation, Prop.~\ref{prop:master}
 & Shaska \cite[Thm.~8.3]{Shaska26} (identical trilinear bracket
   identity, $\Lambda[B,A]+sA[B,\Lambda]+rB[\Lambda,A]=\kappa$, with
   Keller $\iff$ constant); his \S9.2 is titled ``The master equation''.
   Only in v2 --- v1 has the divisorial form
   $\operatorname{Jac}_{u,v}(P,Q)=\kappa\Lambda^{2}$ as Thm.~6.1.
 & Jul 25\\
Parked square: $\det JG=\lambda C^{k}$, order-two vanishing on the
contracted line
 & Shaska \cite[Thm.~7.1]{Shaska26} (= Thm.~6.1 of v1, Jul 22)
 & Jul 22\\
Anchor lemma, Lemma~\ref{lem:anchor}
 & Shaska \cite[Lem.~8.5]{Shaska26}: evaluating the master equation at the
   base point gives $\Lambda(O)\,[B,A](O)=\kappa$, which is our
   $\lambda=C_{0}(1)B_{0}'(1)A_{1}(1)$ after translating $u\mapsto u-1$.
 & Jul 25\\
$S_{3}$ monodromy, non-Galois, discriminant not a square
 & \cite{MO513387} (explicit cubic model, the discriminant, $S_{3}$,
   trivial deck group, and the Campbell/Razar/Wright ``Galois case''
   theorem quoted for
   exactly the reason we quote it); Lou \cite{Lou26}; Mayner
   \cite[\S4.4]{Mayner26}; Shaska \cite[Thm.~4.4(2)]{Shaska26}
 & Jul 20\\
Trace identity (fiber cubic has no subleading term)
 & Mayner \cite[\S4.2]{Mayner26}, in the coordinate $x$ (``the
   $x$-coordinates of the preimages sum to zero''); Shaska
   \cite[Rem.~5.4]{Shaska26}, in the coordinate $s$ (the roots sum to
   $2=\kappa$). Our version is the same phenomenon in the invariant
   coordinate $u$.
 & Jul 21\\
Exact image, Theorem~D
 & ulam.ai \cite[Thm.~4.2]{Ulam26}; Mayner \cite[\S4.3]{Mayner26};
   Shaska \cite[Prop.~5.1]{Shaska26}
 & Jul 20\\
\bottomrule
\end{tabular}
\end{center}

\emph{What is new here.} Two things, plus one reframing.
\begin{enumerate}
\item[(a)] \emph{Theorem~E}, degree minimality in the equivariant class,
with its eight unit-ideal certificates and the degree-$7$ positive
control. We have found nothing comparable; the nearest statements are
compared in Section~\ref{ssec:results}.
\item[(b)] \emph{Theorem~C(i)}, the Hamiltonian lift and the exact
moment-map identity $\nu\circ\Phi=\mu$; and \emph{Lemma~\ref{lem:nogo}},
the no-go lemma for all weights $k\ge1$ and all degrees.
\item[(c)] A \emph{reframing}: Mayner's $\Phi$ is identified as an
explicit $\PC_{n}$-witness through the Adjamagbo--van~den~Essen
formulation \cite{AvdE07}, and tied to $\Psi_{F}$ by
$\gr\Psi_{F}=\Phi^{*}$ (Theorem~C(ii)). The object is not new and the
conclusion $\neg\PC_{n}$ is not new --- it follows formally from
$\neg\DC_{3}$ and is already asserted in secondary sources. What is added
is that the witness is written down, its degrees and momenta computed, and
its non-invertibility checked in two lines from \eqref{eq:collision} with
no filtration or reduction argument.
\end{enumerate}
Minimality certificates for the other weights $k=1$ and $k=3$ are in
preparation and are \emph{not} claimed here.

\emph{Limits of this survey.} The table above rests on a targeted search
of the July--August 2026 record: the arXiv listings for math.AG and
math.RA in that window, the Speyer and Tao threads, the two MathOverflow
questions cited, the ulam.ai note, Lou's derivation, and Mayner's
repository. A negative cannot be proved this way. In particular we were
unable to retrieve MathOverflow question 513390, referenced in
\cite[\S12]{Mayner26}; it returns a 404 and does not appear in the
StackExchange API, so we do not cite it and record its content as
unverified.

\subsection{What is not claimed}\label{ssec:open}
$\JC_{2}$ (the planar Jacobian conjecture), $\DC_{1}$, and $\DC_{2}$
remain open; for $\DC_{1}$ there is a claimed proof by Zheglov
\cite{Zheglov24} currently under review, and see
Bavula--Levandovskyy \cite{BL20} for prior partial results. The
Belov-Kanel--Kontsevich conjecture $\operatorname{Aut}(\W_{n})\cong
\operatorname{Aut}(P_{n})$ \cite{BKK05} is \emph{not} refuted by anything
here: our maps are proper endomorphisms, not automorphisms, so they simply
do not belong to the groups in question. This caution is Mayner's, stated
in the same place as the lift itself (\cite[\S8]{Mayner26}: ``an earlier
draft called this construction `Kontsevich-relevant', which overstates the
connection''), and we repeat it because it is easy to get wrong. The pair
$(\Phi,\Psi_{F})$ is,
however, a sharp stress test for the proposed lifting machinery relating
polynomial symplectomorphisms to Weyl-algebra automorphisms
\cite{KBEY18}, whose status we do not assess here: any such
machinery must be compatible with the existence of a non-invertible
symplectic endomorphism whose canonical quantization is a non-invertible
Weyl endomorphism. All statements are in characteristic zero.

\section{Certified facts about \texorpdfstring{$F$}{F} and its cotangent
lift}\label{sec:facts}

We collect the machine-verified facts used throughout. Script names refer
to Section~\ref{sec:repro}.

\begin{proposition}\label{prop:basic}
For $F$ as in \eqref{eq:themap}:
\begin{enumerate}
\item[(i)] $\det JF=-2$ identically, so $F$ is a Keller map and an
everywhere-local biholomorphism.
\item[(ii)] The triple collision \eqref{eq:collision} holds; the three
source points are distinct points of $\QQ^{3}$. A further rational fiber
point on the generic side: $F(-\tfrac14,-4,48)=(608,-232,1)$.
\item[(iii)] $N:=(JF)^{-1}=\operatorname{adj}(JF)/(-2)$ has polynomial
entries with rational coefficients, of total degree at most $11$.
\end{enumerate}
\emph{(Verified in \texttt{core\_verify.py} and \texttt{weyl\_verify.py};
the nine expanded entries of $N$ are listed in
\texttt{weyl\_endomorphism.txt}.)}
\end{proposition}

Two representative entries of $N$ (columns index the operators $D_{j}$ of
Section~\ref{sec:weyl}):
\[
N_{11}=3x^{6}yz+9x^{5}y^{2}+3x^{5}z+3x^{4}y-4x^{3},\qquad
N_{21}=3x^{4}yz+9x^{3}y^{2}+3x^{3}z+3x^{2}y-3x .
\]

\begin{definition}[Mayner \protect{\cite[\S8]{Mayner26}}]\label{def:phi}
The \emph{cotangent lift} of $F$ is
$\Phi(q,p)=\bigl(F(q),\,(JF(q))^{-\mathsf T}p\bigr)$, a polynomial map
$\CC^{6}\to\CC^{6}$ by Proposition~\ref{prop:basic}(iii). This map, and the
facts recorded in Proposition~\ref{prop:phibasic} that it preserves the
symplectic form exactly, has $\det J\Phi=1$, and is non-injective, are due
to Mayner; our verification is independent. We write
$\Omega=\left(\begin{smallmatrix}0&I_{3}\\-I_{3}&0\end{smallmatrix}\right)$
for the standard symplectic matrix in coordinates
$(x,y,z,p_{1},p_{2},p_{3})$.
\end{definition}

\begin{proposition}[Mayner \protect{\cite[\S8]{Mayner26}}; component
degrees added here]\label{prop:phibasic}
$M:=J\Phi$ satisfies $M^{\mathsf T}\Omega M=\Omega$ identically on
$\CC^{6}$, and $\det M=1$ (as must hold for any symplectic matrix; it is
verified independently). The component degrees of $\Phi$ are
$(7,6,4,9,10,12)$. In particular $\Phi$ has everywhere-invertible
differential, hence open image, hence is dominant.
\emph{(Verified in \texttt{dixmier\_symplectic\_verify.py}; the first two
assertions are Mayner's, verified there by expansion of
$\Phi^{*}\omega-\omega$ and of $\det D\Phi$.)}
\end{proposition}

Since a fiber of $\Phi$ over $(Q,P)$ is
$\{(q,\,JF(q)^{\mathsf T}P):F(q)=Q\}$, the fibers of $\Phi$ are in
canonical bijection with those of $F$: all fiber statements (generic
$3{:}1$, collisions, the missed locus) transfer verbatim between $F$ and
$\Phi$. The three points $a_{i}$ of Theorem~B are exactly the collision
fiber \eqref{eq:collision} decorated with momenta
$p=JF(q_{i})^{\mathsf T}(1,2,3)^{\mathsf T}$.

\section{The Weyl endomorphism and the Dixmier conjecture}
\label{sec:weyl}

Let $\W_{3}=\CC\langle x_{1},x_{2},x_{3},\partial_{1},\partial_{2},
\partial_{3}\rangle$ with $[\partial_{i},x_{j}]=\delta_{ij}$,
$[x_{i},x_{j}]=[\partial_{i},\partial_{j}]=0$, and write
$(x_{1},x_{2},x_{3})=(x,y,z)$. Recall the \emph{order filtration}
$\W_{3}^{(m)}=\operatorname{span}\{x^{\beta}\partial^{\alpha}:
|\alpha|\le m\}$; its associated graded is the commutative polynomial
algebra $\gr\W_{3}\cong\CC[x,y,z,\xi_{1},\xi_{2},\xi_{3}]
=\mathcal{O}(T^{*}\CC^{3})$ via the principal-symbol map, carrying the
standard Poisson bracket. We identify $(\xi_{1},\xi_{2},\xi_{3})$ with the
momenta $(p_{1},p_{2},p_{3})$ of Section~\ref{sec:facts}.

\begin{proposition}[CCR; machine-verified]\label{prop:ccr}
With $D_{j}=\sum_{k}N_{kj}\partial_{k}$:
\begin{enumerate}
\item[(A)] $[D_{j},F_{i}]=\delta_{ij}$ for all $i,j$; equivalently
$JF\cdot N=I_{3}$, nine polynomial identities.
\item[(B)] $[D_{i},D_{j}]=0$ for all $i,j$; equivalently the nine
coefficient polynomials
$\sum_{k}\bigl(N_{ki}\,\partial_{k}N_{lj}-N_{kj}\,\partial_{k}N_{li}\bigr)$
vanish identically.
\end{enumerate}
Hence $x_{i}\mapsto F_{i}$, $\partial_{j}\mapsto D_{j}$ extends to a
well-defined unital endomorphism $\Psi_{F}\colon\W_{3}\to\W_{3}$.
\emph{(Verified by \texttt{weyl\_verify.py}; independently re-checked in
\texttt{dixmier\_symplectic\_verify.py}.)}
\end{proposition}

Identity (A) is a repackaging of $N=(JF)^{-1}$; the substantive content is
(B), the pairwise commutativity of the vector fields $D_{j}$, which
encodes precisely the flatness of the ``inverse coordinate frame'' of the
non-invertible map $F$.

\begin{lemma}\label{lem:gr}
$\Psi_{F}$ preserves the order filtration, and
$\gr\Psi_{F}=\Phi^{*}$ as endomorphisms of
$\mathcal{O}(T^{*}\CC^{3})$.
\end{lemma}

\begin{proof}
$\Psi_{F}(x^{\beta}\partial^{\alpha})=F^{\beta}D^{\alpha}$ has order
$\le|\alpha|$, since each $D_{j}$ has order one and the order of a
product is at most the sum of the orders; hence $\Psi_{F}(\W_{3}^{(m)})\subseteq\W_{3}^{(m)}$ and
the graded endomorphism is defined. On generators,
$\gr\Psi_{F}(x_{i})=F_{i}$ and
$\gr\Psi_{F}(\xi_{j})=\sigma_{1}(D_{j})=\sum_{k}N_{kj}\xi_{k}$. On the
other side, writing $B=(JF)^{-\mathsf T}=N^{\mathsf T}$,
\[
\Phi^{*}(\xi_{j})=\bigl(Bp\bigr)_{j}=\sum_{k}B_{jk}\,\xi_{k}
=\sum_{k}N_{kj}\,\xi_{k},\qquad \Phi^{*}(x_{i})=F_{i}.
\]
The two algebra maps agree on generators, hence everywhere.
\end{proof}

\begin{lemma}\label{lem:strict}
$\ord\Psi_{F}(S)=\ord S$ for every nonzero $S\in\W_{3}$. In particular
every preimage of an order-zero operator has order zero.
\end{lemma}

\begin{proof}
$\Phi$ is dominant (Proposition~\ref{prop:phibasic}), so
$\Phi^{*}$ is injective on $\mathcal{O}(\CC^{6})$. If $\ord S=m$ with
principal symbol $\sigma_{m}(S)\neq0$, then the class of $\Psi_{F}(S)$ in
$\W_{3}^{(m)}/\W_{3}^{(m-1)}$ is
$\Phi^{*}\sigma_{m}(S)\neq0$ by Lemma~\ref{lem:gr}, so
$\ord\Psi_{F}(S)=m$.
\end{proof}

\begin{proof}[Proof of Theorem A]
Well-definedness is Proposition~\ref{prop:ccr}. Injectivity is automatic:
$\W_{3}$ is a simple ring and $\Psi_{F}(1)=1$, so $\ker\Psi_{F}$, a
two-sided ideal not containing $1$, is zero. For non-surjectivity, suppose
$x=\Psi_{F}(S)$. By Lemma~\ref{lem:strict}, $\ord S=0$, i.e.\
$S=g(x,y,z)\in\CC[x,y,z]$, and then $x=g(F_{1},F_{2},F_{3})$. Evaluating
at the first two collision points of \eqref{eq:collision}, which share
their $F$-image but have $x$-coordinates $1$ and $-1$, gives $1=-1$;
contradiction. Hence $x\notin\im\Psi_{F}$.

For $n>3$, write $\W_{n}=\W_{3}\otimes\W_{n-3}$ and set
$\Psi=\Psi_{F}\otimes\id$. It is a unital endomorphism, injective by
simplicity of $\W_{n}$. The order filtration argument goes through
verbatim: $\gr\Psi$ is the pullback along $\Phi\times\id_{\CC^{2(n-3)}}$,
which is dominant, so $\Psi$ is strictly filtered; a preimage of
$x\in\W_{n}$ would give $x=g(F_{1},F_{2},F_{3},x_{4},\dots,x_{n})$, and
evaluation at the two collision points (with the remaining coordinates
fixed) again gives $1=-1$.
\end{proof}

\begin{remark}[Mayner's finer results]\label{rem:mayner}
Much more is true, and was established first in \cite{Mayner26}: the image
of $\Psi_{F}$ is exactly
$\bigoplus_{\alpha\in\mathbb{N}^{3}}\CC[F_{1},F_{2},F_{3}]\,D^{\alpha}$;
none of $x,y,z,\partial_{1},\partial_{2},\partial_{3}$ lies in the image;
the codimension of the image is tabulated in filtration degrees $d\le7$;
and the cokernel is not finitely generated. Note that the restriction of
$\Psi_{F}$ to the order-zero subalgebra
$\CC[x,y,z]\subset\W_{3}$ is precisely the pullback
$F^{*}\colon g\mapsto g\circ F$ --- the classical non-surjective shadow.
The mechanism ``$\Psi_{F}$ automorphism $\Rightarrow$ $F$ invertible'' is
classical; see \cite{vdE00} and \cite[p.~2]{BKK07}. Lemmas
\ref{lem:gr}--\ref{lem:strict} are our packaging of that argument for this
specific $F$.
\end{remark}

\section{The symplectic counterexample and the Poisson conjecture}
\label{sec:poisson}

Equip $\mathcal{O}(\CC^{6})=\CC[x,y,z,p_{1},p_{2},p_{3}]$ with the
standard Poisson bracket
$\{g,h\}=\sum_{i}\bigl(\partial_{p_{i}}g\,\partial_{q_{i}}h-
\partial_{q_{i}}g\,\partial_{p_{i}}h\bigr)$ (sign conventions are
immaterial below). A polynomial map $\Phi\colon\CC^{6}\to\CC^{6}$
satisfies $\{g\circ\Phi,h\circ\Phi\}=\{g,h\}\circ\Phi$ for all $g,h$ if
and only if its Jacobian satisfies $M^{\mathsf T}\Omega M=\Omega$
pointwise; in that case $\Phi^{*}$ is a unital Poisson-algebra
endomorphism of $P_{3}$. The Poisson conjecture $\PC_{n}$, in the
formulation of \cite{AvdE07}, asserts that every such endomorphism of
$P_{n}$ is an automorphism, equivalently that every polynomial map of
$\CC^{2n}$ preserving the standard symplectic form is a polynomial
automorphism.

\begin{proof}[Proof of Theorem B]
The identity $M^{\mathsf T}\Omega M=\Omega$, the value $\det M=1$, the
component degrees, and the collision
$\Phi(a_{1})=\Phi(a_{2})=\Phi(a_{3})=(-\tfrac14,0,0;1,2,3)$ are exact
symbolic facts (Proposition~\ref{prop:phibasic} and
\texttt{dixmier\_symplectic\_verify.py}; the momenta of the $a_{i}$ are
$JF(q_{i})^{\mathsf T}(1,2,3)^{\mathsf T}$). Hence $\Phi^{*}$ is a Poisson
endomorphism of $P_{3}$; it is injective because $\Phi$ is dominant. It is
not surjective: if $x=g\circ\Phi$ for some polynomial $g$, evaluation at
$a_{1},a_{2}$ (equal images, $x$-coordinates $1$ and $-1$) gives $1=-1$.
So $\Phi^{*}$ is a proper endomorphism and $\PC_{3}$ fails. The map $\Phi$
is generically $3{:}1$ because its fibers biject with those of $F$
(Section~\ref{sec:facts}) and $F$ is generically $3{:}1$
(Proposition~\ref{prop:cover} below).

For $n>3$, take $\Phi_{n}=\Phi\times\id_{\CC^{2(n-3)}}$, arranging the
extra coordinates into symplectic pairs; then
$M_{n}^{\mathsf T}\Omega_{2n}M_{n}=\Omega_{2n}$, and non-surjectivity of
$\Phi_{n}^{*}$ follows by the same two-line evaluation. Hence $\PC_{n}$
fails for all $n\ge3$.
\end{proof}

\begin{remark}[Relation to the formal implication, stated exactly]
Given the failure of $\DC_{3}$ (Theorem~A, Mayner's), the failure of
$\PC_{3}$ follows abstractly from $\PC_{3}\Rightarrow\DC_{3}$
\cite{AvdE07}, and is already asserted in secondary sources. The map
$\Phi$ is Mayner's \cite[\S8]{Mayner26}, as are the three properties
verified in Proposition~\ref{prop:phibasic}; his priority section lists
the symplectic lift explicitly. What is added by Theorem~B is the
identification of $\Phi^{*}$ with the Poisson-endomorphism side of
\cite{AvdE07} --- i.e.\ the label $\PC_{n}$ --- together with the
component degrees, the explicit momenta of the collision, and a
non-invertibility proof that is two lines from \eqref{eq:collision}, with
no filtration and no reduction argument. Mayner does not use the words
``Poisson conjecture'' anywhere, and we have found no source that does in
this connection; but the reader should treat Theorem~B as a labelling and
packaging contribution, not as the discovery of a new object.
\end{remark}

\begin{proof}[Proof of Theorem C(ii)]
This is Lemma~\ref{lem:gr}. (Part (i) is proved in
Section~\ref{ssec:moment}.)
\end{proof}

\begin{remark}[Quantization does not restore invertibility]
Theorem C(ii) says the diagram one hopes for --- classical symplectic maps
below, Weyl-algebra maps above, symbol map connecting them --- is realized
exactly here: $\Psi_{F}$ is a filtered quantization of $\Phi$, with
$\gr\Psi_{F}=\Phi^{*}$, and both are proper endomorphisms. Whatever
mechanism produces the failure of injectivity of $F$ survives passage to
the Weyl algebra. This bears on the lifting program of \cite{KBEY18}
(and related work): a correct lifting theory must map the non-invertible
symplectic endomorphism $\Phi$ to a non-invertible object, as the
canonical quantization does. The Belov-Kanel--Kontsevich automorphism
conjecture \cite{BKK05} concerns groups of automorphisms and is untouched
by these examples; see Section~\ref{ssec:open}.
\end{remark}

\section{Equivariant structure: reduction, moment map, cover, image}
\label{sec:equiv}

\subsection{The $\Cstar$-reduction and the master equation}
\label{ssec:reduction}
Give $(x,y,z)$ the weights $(1,-1,-2)$. The invariants of the action are
generated by
\[
u=1+xy,\qquad s=x^{2}z .
\]
The components of $F$ are semi-invariant of weights $(-2,-1,1)$:
reading them against the fiber coordinate $x$,
\begin{equation}\label{eq:collapse}
x^{2}F_{1}=A(u,s),\qquad xF_{2}=B(u,s),\qquad F_{3}=x\,C(u,s),
\end{equation}
with
\begin{equation}\label{eq:ABC}
A=u^{3}s+u(u-1)^{2}(3u+1),\qquad
B=3u^{2}s+9u^{3}-15u^{2}+4u+2,\qquad
C=5-3u-s .
\end{equation}
The invariant combinations of the targets, $p=1+F_{2}F_{3}$ and
$q=F_{1}F_{3}^{2}$, define the \emph{downstairs plane map}
\begin{equation}\label{eq:G}
G(u,s)=\bigl(1+BC,\;AC^{2}\bigr).
\end{equation}

\begin{proposition}[Master equation; \emph{Shaska
\protect{\cite[Thm.~8.3]{Shaska26}}}, obtained independently and
machine-verified here]\label{prop:master}
Let $k\ge1$, $s=x^{k}z$, and let $A,B,C\in\CC[u,s]$ be arbitrary. With
$\Jb(f,g)=f_{u}g_{s}-f_{s}g_{u}$ and $G=(1+BC,\,AC^{k})$,
\begin{equation}\label{eq:master}
\det JG \;=\; C^{k}\,\brk_{k},\qquad
\brk_{k}:=C\,\Jb(B,A)+k\,A\,\Jb(B,C)+B\,\Jb(C,A),
\end{equation}
identically in $\CC[u,s]$. If moreover $F=(A/x^{k},\,B/x,\,xC)$ is a
polynomial map, then $\det JF=-\brk_{k}$ identically. Hence
\[
F\ \text{is Keller}\iff\brk_{k}\in\Cstar .
\]
For the map \eqref{eq:themap} ($k=2$): $\det JG=2C^{2}$, $\brk_{2}=2$, and
$\det JF=-2$.
\end{proposition}

The identity \eqref{eq:master} isolates the design pattern: the Jacobian
of the plane map $G$ ``parks'' a full square $C^{k}$ (for $k=2$, the
square $C^{2}$ on the line $\ell:3u+s=5$), and the fiber scaling
$x\mapsto xC$ cancels it exactly. This divisorial form of the Keller
condition --- $\det JG=\kappa C^{k}$, i.e.\ order-$k$ vanishing on the
contracted locus with constant leading coefficient --- is
\cite[Thm.~7.1]{Shaska26} (Thm.~6.1 of the July 22 version), and the
expanded trilinear identity \eqref{eq:master} itself is
\cite[Thm.~8.3]{Shaska26}; his Remark~8.4 identifies the exponent
$k=\sum_{i\ge2}r_{i}-1$ in general. We obtained \eqref{eq:master}
independently and give the machine verification, but the result is his.
Two further verified facts complete the picture:
\begin{itemize}
\item $G$ \emph{contracts the critical line} $\ell=\{C=0\}$ to the single
point $(p,q)=(1,0)$; restricted to $\ell$,
$(A,B)|_{\ell}=(u^{2}+u,\,4u+2)$ and the one-variable Wronskian
$b\,a'-2a\,b'$ equals $2$ --- the same constant as the Jacobian. The
construction is tight along its critical line.
\item The triple collision \eqref{eq:collision} sits exactly over the
contraction point $(1,0)$: the three fiber points have
$(u,s)$-invariants $u=1$ (the $x=0$ sheet) and $u=-\tfrac12$, $s=\tfrac{13}2$
(twice, the two points on $\{C=0\}$).
\end{itemize}
A heuristic for ``why not $\CC^{2}$'': the mechanism needs a fiber
direction to absorb the parked square; in the plane there is none (and
Moh \cite{Moh83} proved the planar conjecture through degree $100$).

\subsection{The Hamiltonian lift and the moment map (new)}\label{ssec:moment}

\begin{proof}[Proof of Theorem C(i)]
Equivariance of $F$ is the statement \eqref{eq:collapse}:
$F(t\cdot q)=S_{t}F(q)$ with $S_{t}=\operatorname{diag}(t^{-2},t^{-1},t)$.
Differentiating, $JF(t\cdot q)\,T_{t}=S_{t}\,JF(q)$ with
$T_{t}=\operatorname{diag}(t,t^{-1},t^{-2})$, whence
$(JF(t\cdot q))^{-\mathsf T}=S_{t}^{-1}(JF(q))^{-\mathsf T}T_{t}$. Since
the cotangent lift acts on momenta by $T_{t}^{-1}$ upstairs and
$S_{t}^{-1}$ on target momenta, $\Phi$ intertwines the two lifted actions.
The lifted actions are Hamiltonian with the quadratic moment maps
$\mu=xp_{1}-yp_{2}-2zp_{3}$ and $\nu=-2Q_{1}P_{1}-Q_{2}P_{2}+Q_{3}P_{3}$
(weights times $q_{i}p_{i}$), and the exact identity
\[
\nu\circ\Phi=\mu
\]
is verified symbolically in \texttt{dixmier\_symplectic\_verify.py}.
\end{proof}

\begin{remark}[The parked square on the symplectic quotient]
Informally, Theorem C(i) locates the mechanism of
Section~\ref{ssec:reduction} inside symplectic geometry: $\Phi$ descends
to the symplectic quotients of the $\mu$- and $\nu$-level sets, where its
reduced Jacobian degenerates along (the image of) the critical line
$\{C=0\}$ --- the parked square lives on the quotient --- while the group
direction, the fiber coordinate $x$ of \eqref{eq:collapse}, carries the
compensating factor that keeps $\det M\equiv1$ upstairs. We use this only
as an organizing picture; the verified content is the intertwining
identity and the master equation.
\end{remark}

\subsection{Arithmetic of the $3{:}1$ cover}\label{ssec:cover}

\begin{proposition}[machine-verified; \emph{the conclusions of (i) and
(ii) are anticipated}, see the note after the proposition]\label{prop:cover}
Eliminating $s$ from $G=(p,q)$, the minimal polynomial of $u$ over
$\CC(p,q)$ is the cubic
\begin{equation}\label{eq:cubic}
\Delta_{1}\,u^{3}+\bigl((p-1)^{2}-12q\bigr)\,u-4q=0,
\end{equation}
where
\[
\Delta_{1}=p^{3}-4p^{2}-18pq+5p+27q^{2}+34q-2 .
\]
Moreover:
\begin{enumerate}
\item[(i)] (Trace identity.) The cubic has no $u^{2}$ term: the three
preimages of a generic point satisfy
$(1+x_{1}y_{1})+(1+x_{2}y_{2})+(1+x_{3}y_{3})=0$. On the collision fiber:
$1-\tfrac12-\tfrac12=0$.
\item[(ii)] (Discriminant.) $\disc=-4\,\Delta_{1}\Delta_{2}^{2}$ with
$\Delta_{2}=p^{3}-3p^{2}-18pq+3p+54q^{2}+18q-1$, and $\Delta_{1}$ is
irreducible over $\QQ$ --- in particular squarefree of odd degree $3$ in
$p$. Hence $\disc$ is not a square in $\CC(p,q)$: the Galois group of
\eqref{eq:cubic} over $\CC(p,q)$ is the full symmetric group $S_{3}$, and
the cover is not Galois (see Remark~\ref{rem:groundfield} for the descent
to $\CC(p,q)$).
\item[(iii)] ($3{:}1$.) $s$ is a rational function of $(u,p,q)$ (via the
degree-one subresultant, an identity verified exactly), so
$[\CC(u,s):\CC(p,q)]=3$ and $F$ is generically $3$-to-$1$.
\end{enumerate}
\emph{(Verified in \texttt{cover\_verify.py}.)}
\end{proposition}

\begin{remark}[Attribution for Proposition~\ref{prop:cover}]
\label{rem:coverattrib}
Only the coordinate is ours. That the function-field extension has degree
exactly $3$ with full $S_{3}$ monodromy, that the discriminant is not a
square, and that the deck group is therefore trivial, were all established
on July 20, 2026 in \cite{MO513387} (with an explicit cubic model in the
weight-$(-1)$ coordinate $t=y+1/x$, discriminant $-4Q$ with
$Q=27a^{2}c^{2}-18abc+16a+b^{3}c-b^{2}$, and the Campbell/Razar/Wright
Galois-case theorem invoked for precisely the reason we invoke Campbell's
below), and independently by Lou \cite{Lou26} the same day; Mayner
\cite[\S4.4]{Mayner26} has the same $S_{3}$ statement with his
discriminant factorization $-4S^{2}\Delta$, and Shaska
\cite[Thm.~4.4(2)]{Shaska26} states it in the $s$-coordinate with
$\Delta(P,Q)=-4P^{3}+4P^{2}+72PQ-64Q-108Q^{2}$. The trace identity (i) is
the same phenomenon as Mayner's ``the $x$-coordinates of the preimages sum
to zero'' \cite[\S4.2]{Mayner26} and Shaska's ``the three roots always sum
to $2=|\det JF|$'' \cite[Rem.~5.4]{Shaska26}, transported to the invariant
coordinate $u$; the three cubics are genuinely different cubics (different
primitive elements of the same cubic extension), and we do not know
whether the trace-free shape in $u$ specifically is forced.
\end{remark}

\begin{remark}[From the $\QQ$-certificates to $\CC(p,q)$]
\label{rem:groundfield}
The certificates behind Proposition~\ref{prop:cover} operate over $\QQ$,
and irreducibility over $\QQ(p,q)$ does not by itself preclude a
factorization over $\CC(p,q)$; the gap closes as follows. $F$ is an
everywhere-local biholomorphism (Proposition~\ref{prop:basic}(i)), so each
of the three distinct points of the collision fiber \eqref{eq:collision}
has a neighborhood mapped biholomorphically onto a neighborhood of
$(-\tfrac14,0,0)$; every nonempty Zariski-open subset of $\CC^{3}$ meets
the intersection of these image neighborhoods, so some target with
$t\neq0$ and $(p,q)\neq(1,0)$ has at least three $F$-preimages. Its
$F$-fiber is in bijection with the $G$-fiber over $(p,q)$
(Section~\ref{ssec:image}), whose points all satisfy $C\neq0$, where
$\det JG=2C^{2}\neq0$; so $G$ is a local biholomorphism at each of the
$\ge3$ fiber points, and the same openness argument, now downstairs,
produces a Zariski-generic point of $\CC^{2}$ with at least three
$G$-preimages. Since the generic fiber count of a dominant map between
irreducible varieties of equal dimension equals the function-field degree
in characteristic zero, $[\CC(u,s):\CC(p,q)]\ge3$. On the other hand
$s\in\CC(u,p,q)$ by (iii), so $\CC(u,s)=\CC(p,q)(u)$, and $u$ satisfies
the cubic \eqref{eq:cubic}: the degree is $\le3$. Hence the degree is
exactly $3$, the cubic is the minimal polynomial of $u$ over $\CC(p,q)$
--- in particular irreducible over $\CC(p,q)$, with no descent argument
needed --- and the monodromy group is transitive. Finally,
$\disc=-4\Delta_{1}\Delta_{2}^{2}$ is a square in $\CC(p,q)$ if and only
if $-\Delta_{1}$ is; because $\CC[p,q]$ is a UFD, that would force
$-\Delta_{1}$ to be the square of a polynomial, which is impossible:
$\Delta_{1}$ is squarefree --- squarefreeness is the condition
$\gcd(\Delta_{1},\partial_{p}\Delta_{1})=1$, insensitive to extending the
ground field in characteristic zero --- of odd degree $3$ in $p$, while
nonconstant polynomial squares have even degrees. So $\disc$ is not a
square in $\CC(p,q)$, the monodromy is not contained in $A_{3}$, and,
being transitive, it is all of $S_{3}$.
\end{remark}

\begin{remark}[Campbell's theorem makes $S_{3}$ a necessity]
By a theorem of Campbell \cite{Campbell73} --- proved independently, in
general characteristic zero, by Razar and by Wright --- a Keller map whose
extension of rational function fields is Galois is invertible. A
counterexample of covering degree $3$ therefore \emph{had} to be
non-Galois, i.e.\ to have monodromy $S_{3}$ rather than $\mathbb{Z}/3$,
and Proposition~\ref{prop:cover}(ii) confirms that Alp\"oge's map threads
exactly that needle. This observation is also not ours: it is made in
\cite{MO513387} on July 20, 2026, with the same citation. Whether the
trace identity (i) is forced by the mechanism or is extra tuning is an
open question worth testing on any future examples with other weights.
\end{remark}

\subsection{The image theorem}\label{ssec:image}

The statement of Theorem~D is not new; see the priority note in its
statement and Section~\ref{ssec:credit}. The proof below is our own and is
included because the $G$-fiber correspondence it sets up is reused in
Theorem~E, and because the transfer $\im\Phi=\CC^{6}\setminus(\Gamma\times
\CC^{3})$ needs it. Readers who prefer the published treatments should
consult \cite[Thm.~4.2]{Ulam26}, which routes the fiber count through a
binary cubic $2pS^{3}-qS^{2}T+2ST^{2}-rT^{3}$ and an incidence variety in
$\CC^{3}\times\mathbb{P}^{1}$, or \cite[Prop.~5.1]{Shaska26}, which routes
it through the contracted line.

\begin{proof}[Proof of Theorem D]
Write the target coordinates as $(X,Y,t)$.

\emph{The plane $t=0$ is covered.} On the sheet $x=0$ the map is
$F(0,y,z)=(z+4y^{2},\,y,\,0)$, visibly bijective onto the plane
$\{t=0\}$.

\emph{Targets with $t\neq0$.} By \eqref{eq:collapse}, a preimage with
$F_{3}=t\neq0$ has $x\ne0$ and $C(u,s)\neq0$, and corresponds to a point
of the $G$-fiber over
\[
(p,q)=(1+tY,\;t^{2}X),
\]
with $x=t/C(u,s)$ recovered uniquely; conversely every $G$-fiber point
with $C\neq0$ arises this way. If $(p,q)\ne(1,0)$, every $G$-fiber point
automatically has $C\neq0$ (since $C=0$ forces $(p,q)=(1,0)$), so the
$F$-fiber over $(X,Y,t)$ is in bijection with the $G$-fiber over $(p,q)$.
The axis targets $(0,0,t)$, $t\neq0$ --- exactly the case $(p,q)=(1,0)$
--- are hit exactly once, by $(t/2,0,0)$: the system $A=B=0$, $C\neq0$
has the single solution $(u,s)=(1,0)$, where $C=2$ (reduced Gr\"obner
basis over $\QQ$, asserted in \texttt{min\_verify.py} part
\texttt{axis}), whence $x=t/2$, $y=0$, $z=0$.

\emph{The line $q=0$ is covered.} On the $u=0$ sheet $q\equiv0$ and
$p=11-2s$ takes every value; so every $(p,q)$ with $q=0$ has a $G$-fiber
point, and roots $u_{0}=0$ of the cubic --- which occur only when the
constant term $-4q$ vanishes --- never need to be lifted.

\emph{From cubic roots to fiber points.} Fix a target $(p,q)$ with
$q\neq0$ and let $u_{0}$ be any root of the cubic \eqref{eq:cubic}. Viewed
as polynomials in $s$, $p-P$ and $q-Q$ have degrees $2$ and $3$ with
$s$-leading coefficients $-3u^{2}$ and $u^{3}$ respectively (from
$BC$ and $AC^{2}$), and the cubic \eqref{eq:cubic} is an irreducible
factor of their $s$-resultant --- indeed the resultant equals $-u^{3}$
times \eqref{eq:cubic}, the extraneous factor being supported on
$\{u=0\}$ (\texttt{cover\_verify.py} extracts \eqref{eq:cubic} from
precisely this resultant). Since $q\neq0$, the constant term $-4q$ of the
cubic is nonzero, so $u_{0}\neq0$ and both $s$-leading coefficients
$-3u_{0}^{2}$, $u_{0}^{3}$ are nonzero; hence specializing $u=u_{0}$
commutes with taking the resultant, and the specialized resultant
vanishes because the cubic does. Two univariate polynomials with
nonvanishing leading coefficients and vanishing resultant have a common
root: there is $s_{0}\in\CC$ with $(u_{0},s_{0})\in G^{-1}(p,q)$. So the
$G$-fiber over any $(p,q)$ with $q\neq0$ is nonempty whenever the cubic
has a root.

\emph{Emptiness locus of $G$.} By the previous two steps, the fiber can
be empty only where the cubic has no root at all: only where its leading
coefficient $\Delta_{1}$ and its linear coefficient $(p-1)^{2}-12q$ both
vanish while the constant term $-4q$ does not. Solving, the unique such
point is $(p,q)=(\tfrac73,\tfrac{4}{27})$, and a Gr\"obner certificate
over $\QQ$ shows the ideal generated by the $G$-equations at that point
is the \emph{unit ideal}: that fiber is indeed empty (Nullstellensatz;
\texttt{min\_verify.py}, \texttt{cover\_verify.py}).

\emph{The missed curve.} Pulling $(p,q)=(\tfrac73,\tfrac{4}{27})$ back
through $(p,q)=(1+tY,t^{2}X)$ gives, for each $t\neq0$, the unique missed
target $\bigl(\tfrac{4}{27t^{2}},\tfrac{4}{3t},t\bigr)$, i.e.\ the curve
$\Gamma$. Every other target is attained, so
$\im F=\CC^{3}\setminus\Gamma$, and
$\im\Phi=\CC^{6}\setminus(\Gamma\times\CC^{3})$ by the fiber
correspondence of Section~\ref{sec:facts}.

\emph{Fiber counts.} Generic fibers have three points. Over
$\{\Delta_{2}=0\}$ (with $\Delta_{1}\neq0$) the fiber still has three
\emph{distinct} points: $\det JG=2C^{2}$ vanishes only on $\{C=0\}$, which
$G$ contracts to the single point $(1,0)$, so no two of the three sheets
can genuinely coincide over any $(p,q)\neq(1,0)$; the vanishing of
$\disc=-4\Delta_{1}\Delta_{2}^{2}$ on $\{\Delta_{2}=0\}$ is therefore
\emph{apparent}, the invariant $u=1+xy$ failing to separate two unramified
sheets, whence $\Delta_{2}$ enters squared. Over
$\{\Delta_{1}=0\}\setminus\{(\tfrac73,\tfrac{4}{27})\}$ the cubic
\eqref{eq:cubic} drops degree $3\to1$ --- its $u^{2}$ term being absent by
the trace identity, two sheets escape to infinity at once --- and exactly
one preimage remains; at $(\tfrac73,\tfrac{4}{27})$, i.e.\ over $\Gamma$,
all three sheets have escaped and the fiber is empty. The achievable
set-theoretic fiber sizes are therefore exactly $\{3,1,0\}$.
\end{proof}

So $F$ is an everywhere-local biholomorphism of $\CC^{3}$, generically
$3{:}1$, surjective except for one explicit punctured rational curve --- a
picture worth keeping: over $\Gamma$ the missing preimages have escaped
to infinity along $\{\Delta_{1}=0\}$.

\begin{remark}[Transport to $\Psi_{F}$]\label{rem:transport}
At order zero, the deficiency of $\Psi_{F}$ \emph{is} the deficiency of
$F^{*}$: $\CC[F]\subsetneq\CC[x,y,z]$, an inclusion of polynomial algebras
dual to $F$ itself, whose geometric failure of surjectivity is exactly the
curve $\Gamma$. In higher filtration degrees, Mayner's computations
\cite{Mayner26} (exact image
$\bigoplus_{\alpha}\CC[F]D^{\alpha}$, codimension table for $d\le7$,
non-finitely-generated cokernel) quantify the deficiency of $\Psi_{F}$; at
the symbol level Lemma~\ref{lem:gr} bounds the image by
$\im\Phi^{*}$, whose geometric deficiency is $\Gamma\times\CC^{3}$
(Theorem D). We expect the growth of $\operatorname{coker}\Psi_{F}$ to be
governed by functions along $\Gamma$ and the $S_{3}$ monodromy of
Proposition~\ref{prop:cover}, but we leave a precise statement to future
work.
\end{remark}

\section{Degree minimality in the equivariant class}\label{sec:min}

Fix the weights $(1,-1,-2)$, i.e.\ $k=2$, $u=1+xy$, $s=x^{2}z$, and
consider all maps of the shape \eqref{eq:collapse}:
$F=(A/x^{2},B/x,xC)$ with $A,B,C\in\CC[u,s]$. Writing
$A=A_{0}(u)+A_{1}(u)s+A_{2}(u)s^{2}+\cdots$ and likewise for $B,C$, the
condition that $F$ be a polynomial map (\emph{lift conditions}) is
\[
(u-1)^{2}\mid A_{0},\qquad B_{0}(1)=0,
\]
and by Proposition~\ref{prop:master} the Keller condition is
$\brk_{2}=\lambda\in\Cstar$. We refer to this family as \emph{the
equivariant class}; the map \eqref{eq:themap} is the member
\eqref{eq:ABC} with $\lambda=2$.

\subsection{The degree dictionary and normalization}
The monomial substitution $u^{i}s^{j}\mapsto x^{i+2j}y^{i}z^{j}$ is
injective on exponent pairs, so there is no cancellation of leading terms
and degrees upstairs are read off exactly. The condition $\deg F\le6$
caps the $s$-expansion coefficients precisely as follows:
\[
\begin{array}{l|ccc}
 & s^{0}\text{-part} & s^{1}\text{-part} & s^{2}\text{-part}\\\hline
A: & \deg\le4 & \deg\le2 & \deg\le1\\
B: & \deg\le3 & \deg\le2 & \text{constant}\\
C: & \deg\le2 & \deg\le1 & \text{---}
\end{array}
\]
Conjugating by the two diagonal $\Cstar$-automorphisms upstairs scales
$(A,B,C)\mapsto(\nu r^{-2}A,\,\nu r^{-1}B,\,\nu rC)$ and
$\lambda\mapsto\nu^{3}r^{-2}\lambda$; this legitimately normalizes
$\lambda=1$ together with one further nonzero leading coefficient to $1$
in each case below.

\subsection{The anchor lemma}

\begin{lemma}[Anchor; \emph{Shaska \protect{\cite[Lem.~8.5]{Shaska26}}},
obtained independently here]\label{lem:anchor}
For every Keller map in the equivariant class, of any degree,
\[
\lambda=C_{0}(1)\cdot B_{0}'(1)\cdot A_{1}(1).
\]
In particular every Keller lift --- counterexample or automorphism ---
satisfies $C_{0}(1)\neq0$, $B_{0}'(1)\neq0$, and $A_{1}(1)\neq0$.
\end{lemma}

\begin{proof}
Evaluate $\brk_{2}=C\,\Jb(B,A)+2A\,\Jb(B,C)+B\,\Jb(C,A)$ at
$(u,s)=(1,0)$. The lift conditions give $A(1,0)=A_{0}(1)=0$ and
$A_{u}(1,0)=A_{0}'(1)=0$ (double root), and $B(1,0)=B_{0}(1)=0$; so the
second and third terms vanish, and in the first,
$\Jb(B,A)(1,0)=B_{u}A_{s}-B_{s}A_{u}=B_{0}'(1)A_{1}(1)$. Multiplying by
$C(1,0)=C_{0}(1)$ gives the identity. (Machine-verified at full
$\deg\le6$ generality in \texttt{min\_verify.py}; the derivation above is
degree-free.) For the map \eqref{eq:ABC}: $2=2\cdot1\cdot1$.
\end{proof}

\begin{remark}[Attribution for Lemma~\ref{lem:anchor}]
Shaska states this as: evaluating the master equation at the base point
$O$ of the quotient plane gives $\Lambda(O)\,\Jb(B,A)(O)=\kappa$, whence
$\Lambda(O)\neq0$ and $dA\wedge dB|_{O}\neq0$
\cite[Lem.~8.5]{Shaska26}. Translating his base point to $(u,s)=(1,0)$ and
using the lift conditions to kill $A_{u}(1,0)$ turns
$\Jb(B,A)(O)=B_{u}A_{s}-B_{s}A_{u}$ into $B_{0}'(1)A_{1}(1)$, which is the
displayed factorization: the two statements are the same lemma, and his is
earlier (July 25, 2026, in v2 of \cite{Shaska26}). We keep our formulation
because the factored form is what the case analysis of Theorem~E consumes.
\end{remark}

\subsection{The no-go lemma}

The following appears to be new; we have found no analogue in
\cite{Mayner26,Shaska26,Ulam26,Lou26,MO513387}. The nearest published
statements are \cite[Rem.~8.8]{Shaska26}, which handles the case
$\Lambda$ constant (where the graded problem in dimension $n$ collapses to
the ordinary Jacobian conjecture in dimension $n-1$), and
\cite[Thm.~10.10]{Shaska26}, which handles $A$ and $B$ both of degree one
in the invariants. Lemma~\ref{lem:nogo} constrains the $s$-degree rather
than the total degree, and is what makes Branch~I of Theorem~E finite.

\begin{lemma}[No-go; all weights $k\ge1$; \emph{new}]\label{lem:nogo}
Fix $k\ge1$ and weights $(1,-1,-k)$. If $C$ is $s$-free and $A,B$ are at
most affine in $s$, then every Keller lift in the class is a polynomial
automorphism of $\CC^{3}$ --- at any degree.
\end{lemma}

\begin{proof}
Write $A=A_{0}(u)+A_{1}(u)s$, $B=B_{0}(u)+B_{1}(u)s$, $C=C_{0}(u)$, and
recall the invariant outputs $p-1=BC$ and $q=AC^{k}$. Note first that
$C_{0}\neq0$ as a polynomial (else $\brk_{k}=0$), and that if
$A_{1}=B_{1}=0$ then $\brk_{k}=0$ and the map is not Keller; so at least
one of $A_{1},B_{1}$ is nonzero. Throughout, $(X,Y,T)$ denote the target
coordinates of $F$, so that $p=1+YT$ and $q=XT^{k}$ are rational in the
target.

\emph{Main case $A_{1}B_{1}\neq0$.} The bracket $\brk_{k}$ is affine in
$s$; setting its $s$-coefficient to zero and dividing by
$A_{1}B_{1}C_{0}$ gives
$B_{1}'/B_{1}-A_{1}'/A_{1}-(k-1)\,C_{0}'/C_{0}=0$, i.e.\
$B_{1}=\kappa A_{1}C_{0}^{\,k-1}$ for a constant $\kappa\in\Cstar$ (a
logarithmic-derivative computation). With this substitution the bracket
collapses to the verified identity
\[
\brk_{k}=A_{1}\cdot\bigl(C_{0}B_{0}-\kappa\,C_{0}^{\,k}A_{0}\bigr)' .
\]
A polynomial times a derivative being the nonzero constant $\lambda$
forces $A_{1}\in\Cstar$ and
$C_{0}B_{0}-\kappa C_{0}^{k}A_{0}=\tilde\lambda u+\mu$ with
$\tilde\lambda=\lambda/A_{1}\neq0$. A second verified identity says
that, with these forced shapes, the $s$-terms of $(p-1)-\kappa q$ cancel
\emph{identically}:
\[
(p-1)-\kappa q\;=\;C_{0}B_{0}-\kappa\,C_{0}^{\,k}A_{0}
\;=\;\tilde\lambda u+\mu .
\]
Now invert $F$ rationally, in three steps. First,
$u=\bigl((p-1)-\kappa q-\mu\bigr)/\tilde\lambda$ is a polynomial in
$(p,q)$, hence rational in the target. Second, at fixed $u$ the relation
$p-1=(B_{0}+B_{1}s)\,C_{0}$ is affine in $s$ with $s$-coefficient
$B_{1}C_{0}=\kappa A_{1}C_{0}^{\,k}\neq0$ in $\CC(u)$, so
$s=\bigl((p-1)-B_{0}C_{0}\bigr)/\bigl(\kappa A_{1}C_{0}^{\,k}\bigr)$ is
rational in the target. Third, $x=F_{3}/C(u,s)=T/C_{0}(u)$, then
$y=(u-1)/x$ and $z=s/x^{k}$, are rational in the target. So $F$ is
birational, and a birational Keller map is a polynomial automorphism
\cite[Corollary~1.1.35]{vdE00} (equivalently
\cite[Theorem~2.1]{BCW82}; the birational case goes back to Keller
\cite{Keller39}).

\emph{Edge case $B_{1}=0$, $A_{1}\neq0$.} Here the bracket collapses to
$\brk_{k}=A_{1}\,(C_{0}B_{0})'$ --- the $\kappa=0$ instance of the
displayed identity --- so again $A_{1}\in\Cstar$ and
$C_{0}B_{0}=\tilde\lambda u+\mu$ with $\tilde\lambda\neq0$. Since $B$ is
$s$-free, $p-1=B_{0}C_{0}=\tilde\lambda u+\mu$, so $u$ is affine in $p$;
then $q=(A_{0}+A_{1}s)\,C_{0}^{\,k}$ gives
$s=\bigl(q\,C_{0}^{-k}-A_{0}\bigr)/A_{1}$, rational in the target, and
$x,y,z$ are recovered as above: $F$ is birational, hence an
automorphism.

\emph{Edge case $A_{1}=0$, $B_{1}\neq0$.} Here the bracket reduces to
the relation $B_{1}\,(C_{0}^{k}A_{0})'=-\lambda\,C_{0}^{\,k-1}$. Compare
orders of vanishing at $u=1$, writing $a=\ord_{1}A_{0}$ and
$c=\ord_{1}C_{0}$: the lift condition $(u-1)^{k}\mid A_{0}$ gives
$a\ge k$, so the left side has
$\ord_{1}=\ord_{1}B_{1}+kc+a-1$ (the derivative drops the order by
exactly one in characteristic zero), while the right side has
$\ord_{1}=(k-1)c$; equality forces $\ord_{1}B_{1}+c+a=1$, impossible for
$k\ge2$ since $a\ge k\ge2$ --- this subcase is empty, consistent with
Theorem~E. For $k=1$ the relation reads $B_{1}\,(C_{0}A_{0})'=-\lambda$:
a product of polynomials equal to a nonzero constant, so
$B_{1}\in\Cstar$ and $C_{0}A_{0}$ is affine in $u$ with nonzero slope.
Then $q=AC=A_{0}C_{0}$ is affine in $u$ with nonzero slope: $u$ is
affine in $q$; then $p-1=(B_{0}+B_{1}s)C_{0}$ recovers $s$ rationally
($B_{1}C_{0}\neq0$), and $x,y,z$ follow as above: $F$ is birational,
hence an automorphism. \emph{(The two identities displayed in the main
case are verified in \texttt{core\_verify.py}.)}
\end{proof}

\begin{corollary}\label{cor:exit}
A counterexample in the class must either put $s$ inside $C$, or put
$s^{2}$-terms inside $A$ or $B$. Alp\"oge's map takes the first exit:
$C=5-3u-s$.
\end{corollary}

\subsection{The minimality theorem}

\begin{proof}[Proof of Theorem E]
By the degree dictionary, $\deg F\le6$ means the caps in the table above.
We split on the shape of $C$.

\emph{Branch I: $C$ is $s$-free} ($C=C_{0}(u)$, $\deg C_{0}\le2$).
\begin{itemize}
\item If $C_{0}$ is constant, then $G$ is an affine map composed with a
planar Keller pair of degree $\le4$, which is invertible by Moh's theorem
\cite{Moh83}; the lift is an automorphism.
\item If $A,B$ are at most affine in $s$, Lemma~\ref{lem:nogo} applies:
the lift is an automorphism.
\item There remains the \emph{$s^{2}$-sector}: $C_{0}$ nonconstant and
$(A_{2},B_{2})\neq(0,0)$. After normalization this is six polynomial
systems in the coefficient unknowns (nonvanishing of the leading
coefficient of $C_{0}$ imposed by a Rabinowitsch variable). For all six,
msolve returns the \emph{unit ideal} over $\QQ$: by the Nullstellensatz
the sector is empty --- it contains no Keller map whatsoever.
\end{itemize}

\emph{Branch II: $C_{1}\neq0$} ($C$ genuinely involves $s$). Two systems
at \emph{full} generality --- all coefficients of $A$ (through $s^{2}$,
degrees capped), $B$ (through $s^{2}$), $C_{0}$, and $C_{1}$ free, with
only the legitimate normalizations $\lambda=1$ and $C_{1}$ monic
(respectively constant $=1$); no hand-derived reduction is load-bearing.
Both return the unit ideal over $\QQ$: nothing exists in this sector
below degree $7$, automorphisms included.

The eight systems are:
\begin{center}
\begin{tabular}{lllll}
\toprule
Leaf & Sector & Equations & Unknowns & Result ($\QQ$ and
$\mathbb{F}_{32003}$)\\
\midrule
II-f1 & $C_{1}$ monic linear & 26 & 19 & $(1)$\\
II-f0 & $C_{1}=1$ & 23 & 18 & $(1)$\\
I-B2-g2 & $B_{2}=1$, $\deg C_{0}=2$ & 23 & 18 & $(1)$\\
I-B2-g1 & $B_{2}=1$, $\deg C_{0}=1$ & 19 & 17 & $(1)$\\
I-A2L-g2 & $A_{2}$ monic linear, $\deg C_{0}=2$ & 20 & 17 & $(1)$\\
I-A2L-g1 & $A_{2}$ monic linear, $\deg C_{0}=1$ & 17 & 16 & $(1)$\\
I-A2c-g2 & $A_{2}=1$, $\deg C_{0}=2$ & 19 & 16 & $(1)$\\
I-A2c-g1 & $A_{2}=1$, $\deg C_{0}=1$ & 16 & 15 & $(1)$\\
\bottomrule
\end{tabular}
\end{center}

\emph{Positive control.} At the degree-$7$ caps, the torus-scaled map
\eqref{eq:ABC} satisfies the identical normalized system ($\brk_{2}=1$,
$C_{1}$-coefficient $=1$) --- verified exactly by substitution: with
$\nu=2^{-1/5}$ (exact, $\nu^{5}=\tfrac12$) and $r=-1/\nu$, the scaled
triple $(\nu r^{-2}A,\,\nu r^{-1}B,\,\nu rC)$ has $C_{1}$-coefficient
$1$ and bracket identically $1$; the asserted form of this check is
\texttt{min\_verify.py} part \texttt{d7control}. (An msolve reduced-basis
run on the degree-$7$ control system did not terminate within a
$600$-second cap --- expected for a nonempty positive-dimensional
variety; the exact substitution above is the control, and the
computational record is
\texttt{certificates/D7CONTROL-NEGATIVE-RESULT.md} in the repository.)
The pipeline provably contains the counterexample at degree $7$ and
provably contains nothing at degree $\le6$.

All eight emptiness certificates were computed over $\QQ$ (msolve
multi-modular F4) and again mod $32003$. The cross-system reproduction
reaches the six Branch-I leaves: SymPy's Gr\"obner engine returns the
unit ideal for all six in about a second (\texttt{min\_verify.py}
part \texttt{I}). It does not reach the two Branch-II leaves: part
\texttt{II} exhausts its $150$-second per-leaf alarm on both on the machines
we have run it on, as does
the longer $240$-second attempt over $\QQ$, so SymPy returns no verdict
there and Branch-II emptiness rests on the stored msolve certificates
\texttt{out\_II-*\_q.txt} and \texttt{out\_II-*\_p.txt}, each the
reduced basis $[1]$. SymPy additionally verifies the structural layer
identities (Lemma~\ref{lem:anchor} and the layer factorizations) used to
\emph{organize}, but not to prove, the case split.
\end{proof}

\begin{remark}[Anatomy of the exit at degree $7$]
In the sector $A_{2}=B_{2}=0$, $C_{1}\neq0$, a layer computation shows
the top $s$-layer of the bracket forces $B_{1}^{3}=\rho\,A_{1}^{2}C_{1}$
for a constant $\rho$, i.e.\
$3\deg B_{1}=2\deg A_{1}+\deg C_{1}$. Under the $\deg\le6$ caps only the
patterns $(0,0,0)$ and $(1,1,1)$ survive, and both die against
Lemma~\ref{lem:anchor}; the example lives on the pattern
$(\deg B_{1},\deg A_{1},\deg C_{1})=(2,3,0)$ --- $6=6+0$ --- which forces
$\deg F_{1}=7$. Degree $7$ is not an accident; it is the first rung the
mechanism can reach.

This layer relation is a special case of a published one. Shaska
\cite[Thm.~10.6]{Shaska26} shows that the leading forms
$\hat A,\hat B,\hat\Lambda$ of any graded Keller map are, up to
constants, products of powers of a common set of linear forms, subject to
the Diophantine system
$\beta_{i}+r\gamma_{i}=(\nu_{B}/e)m_{i}$,
$\alpha_{i}+s\gamma_{i}=(\nu_{A}/e)m_{i}$; and his Example~10.9 works out
that system at the degrees $\mathbf d=(4,3,1)$ of Alp\"oge's map, finding
$h=u(3u+v)$, $m=(1,1)$, $\beta=(2,1)$, $\gamma=(0,1)$, $\alpha=(3,1)$ and
concluding that ``the leading data of the announced counterexample is thus
the solution of a small system of linear equations in nonnegative
integers, and not an accident''. Our relation
$B_{1}^{3}=\rho A_{1}^{2}C_{1}$ is the same constraint read one $s$-layer
at a time, and the conclusion above is the same conclusion; his is earlier
(July 25, 2026) and more general. What the remark adds is the pairing of
that constraint with the anchor identity to \emph{exclude} the surviving
low-degree patterns --- which is exactly the content that
Theorem~E certifies and that \cite[Rem.~10.8]{Shaska26} explicitly says
the stratification alone does not give (``emptiness must therefore be
proved stratum by stratum'').
\end{remark}

\begin{remark}[Scope caveats, stated plainly]\label{rem:caveats}
(i) Theorem E is relative to this symmetry class: weights $(1,-1,-2)$ and
the lift shape \eqref{eq:collapse}. Other weights $(1,-1,-k)$ and
non-equivariant maps are untouched; ``no degree-$\le6$ counterexample in
$\CC^{3}$'' remains open --- unconditionally, invertibility of Keller
maps is known only through degree $2$, in every dimension, by Wang's
theorem \cite{Wang80}. (ii) The characteristic-zero emptiness
certificates rest on msolve's multi-modular F4 over $\QQ$; the mod-$p$
runs concur. An independent certificate checker that replays the
unit-ideal cofactor identities has not yet been built, so a reader who
does not trust msolve over $\QQ$ has at present only the mod-$32003$
cross-reproduction, and that in a second system for the six Branch-I
leaves alone. (iii) Moh's theorem and Campbell's theorem are quoted
from the literature. (iv) Certificates for the weights $k=1$ and $k=3$ are
in preparation and are not claimed here. (v) Emptiness \emph{inside} a
class is not emptiness outside it: ulam.ai \cite[Thm.~5.1,
Cor.~5.3]{Ulam26} produces a $(\lambda,a,c,H)$-family with the same cubic
geometry, and nonproper Keller maps of every generic degree $d\ge3$; Gao
\cite{Gao26} produces further explicit counterexamples in every dimension
$>2$ by a tangent-sweep construction; Migus \cite{Migus26} determines the
possible generic degrees over $\mathbb{R}$. None of these is a degree-$\le6$
member of our equivariant class, but they are a reminder that the map is
not isolated in any broad sense --- a point Mayner records as a correction
to his own rigidity claim \cite[\S12]{Mayner26}.
\end{remark}

\begin{remark}[Stable reductions]
By Bass--Connell--Wright \cite{BCW82}, Yagzhev \cite{Yagzhev80}, and
Dru\.zkowski \cite{Druzkowski83}, the example yields cubic-homogeneous
(even cubic-linear-form) counterexamples in some higher $\CC^{N}$; in the discussion thread of
\cite{Speyer26} an explicit cubic-homogeneous reduction in $24$ variables
was described (Thompson), and Long \cite{Long26} tracks a conservative
Bass--Connell--Wright reduction of the announced map to $79$ variables en
route to the Gaussian Moments conjecture. Explicit degree/dimension
bookkeeping, and the smallest such $N$, would be a worthwhile follow-up,
as would the minimality question for the other weights $(1,-1,-k)$ --- the
master equation \eqref{eq:master} is already general, and the same
certificate pipeline applies.
\end{remark}

\begin{remark}[Other consequences drawn elsewhere]
For orientation, the July 2026 consequence-drawing around this map
includes: the Dixmier conjecture, Mayner \cite{Mayner26}; the Gaussian
Moments conjecture, Long \cite{Long26}; the Hessian conjecture in five
variables, Meng--Yang \cite{MengYang26}; the separable Jacobian conjecture
in characteristic $2$, Huq-Kuruvilla \cite{HuqKuruvilla26}; quantum
inequivalence of scalar field redefinitions, Zhu \cite{Zhu26}; the
structure of the space of Keller maps, Jelonek \cite{Jelonek26}; real
generic degrees, Migus \cite{Migus26}; further families, ulam.ai
\cite{Ulam26} and Gao \cite{Gao26}; and the graded classification, Shaska
\cite{Shaska26}. The present note contributes to none of those lines
except the last.
\end{remark}

\section{Reproducibility}\label{sec:repro}

All artifacts sit beside this note. Trust model, stated exactly (it
matches Remark~\ref{rem:caveats}(ii)): every polynomial identity in this
note is an exact symbolic check asserted in SymPy 1.14
(arbitrary-precision rational arithmetic). The characteristic-zero
emptiness certificates of Theorem~E rest on msolve 0.10.1's certified
multi-modular F4 over $\QQ$; SymPy independently reproduces the
mod-$32003$ runs for the six Branch-I leaves, so the cross-system
reproduction concerns those mod-$p$ computations, while the two
Branch-II leaves and all eight characteristic-zero certificates rest on
msolve alone. No numerical approximation occurs anywhere. A
self-contained independent certificate checker (replaying the unit-ideal
cofactor identities) is desirable and has not yet been built.

\begin{itemize}
\item \texttt{core\_verify.py} --- the map \eqref{eq:themap}, $\det JF=-2$,
the collisions, the collapse \eqref{eq:collapse}, the master equation
\eqref{eq:master} (all $k$, with symbolic function coefficients), the
critical-line facts, and the two identities of Lemma~\ref{lem:nogo}.
\item \texttt{cover\_verify.py} --- the cubic \eqref{eq:cubic}, trace
identity, discriminant factorization, $S_{3}$ monodromy, rationality of
$s$, and the image/escape analysis of Theorem D.
\item \texttt{weyl\_verify.py} --- polynomiality of $N$, the CCR
identities (A) and (B) of Proposition~\ref{prop:ccr}; writes
\texttt{weyl\_endomorphism.txt}, the nine fully expanded operator
coefficients.
\item \texttt{dixmier\_symplectic\_verify.py} --- $M^{\mathsf T}\Omega
M=\Omega$, $\det M=1$, component degrees, the $\CC^{6}$ triple collision,
the moment-map identity $\nu\circ\Phi=\mu$, and an independent re-check
of the CCR.
\item \texttt{min\_verify.py} --- Lemma~\ref{lem:anchor} at full
$\deg\le6$ generality, the layer identities, the image certificate at
$(\tfrac73,\tfrac{4}{27})$ over $\QQ$ (part \texttt{ident}); a SymPy
Gr\"obner harness for the eight leaf systems mod $32003$ (parts
\texttt{I}, \texttt{II}); the degree-$7$ positive control by exact
scaled substitution (part \texttt{d7control}); the upstairs identity
$\det JF=-\brk_{k}$ for $k=1,2,3$ with free symbolic coefficients (part
\texttt{kdet}); and the axis-target uniqueness certificate (part
\texttt{axis}).
\item \texttt{ms\_*\_c0.ms}, \texttt{ms\_*\_c32003.ms},
\texttt{out\_*.txt} --- msolve input files for the eight leaf systems in
characteristic $0$ and mod $32003$, and the sixteen outputs
(\texttt{out\_*\_q.txt} = characteristic $0$,
\texttt{out\_*\_p.txt} = mod $32003$); every output is the reduced
Gr\"obner basis $[1]$. The scripts that generated these inputs are not
part of the deposit.
\end{itemize}

\begin{verbatim}
python3 -m venv venv && venv/bin/pip install sympy
venv/bin/python core_verify.py
venv/bin/python cover_verify.py
venv/bin/python weyl_verify.py
venv/bin/python dixmier_symplectic_verify.py
venv/bin/python min_verify.py ident
venv/bin/python min_verify.py d7control
venv/bin/python min_verify.py kdet
venv/bin/python min_verify.py axis
venv/bin/python min_verify.py I    # SymPy GB, six Branch-I leaves, mod 32003
venv/bin/python min_verify.py II   # SymPy GB, two Branch-II leaves, mod 32003
brew install msolve
for f in ms_I-*_c0.ms ms_II-*_c0.ms; do b=${f#ms_}; \
  msolve -g 2 -f "$f" -o "out_${b%_c0.ms}_q.txt"; done          # char 0
for f in ms_I-*_c32003.ms ms_II-*_c32003.ms; do b=${f#ms_}; \
  msolve -g 2 -f "$f" -o "out_${b%_c32003.ms}_p.txt"; done      # mod 32003
\end{verbatim}
Expected: every \texttt{out\_*\_q.txt} and \texttt{out\_*\_p.txt} contains
the reduced Gr\"obner basis $[1]$, matching the stored outputs. The msolve
runs above are not expected to time out; \texttt{min\_verify.py} part
\texttt{II} is expected to print \texttt{TIMEOUT} on both leaves; a
timeout is inconclusive, not a failure, and the Branch-II proof is the
stored msolve output.

\section*{Acknowledgments}
We are grateful to W.~G.~P.~Mayner, whose note and report \cite{Mayner26}
first put both the Weyl-algebra endomorphism \emph{and} the symplectic
cotangent lift on public record, and whose computations go beyond ours at
several points; to T.~Shaska \cite{Shaska26}, whose graded framework
contains, in greater generality and earlier, the master equation, the
base-point identity, the divisorial form of the Keller condition, and the
leading-form analysis that we had arrived at independently; to the
anonymous author of \cite{Ulam26} for the earliest published fiber-count
and image theorem; and to the participants of the Secret Blogging Seminar
discussion \cite{Speyer26} --- in particular the explicit $24$-variable
cubic-homogeneous reduction described there (Thompson). Discovering that a
result one has proved was proved first by someone else is the ordinary
condition of working on a problem that is three weeks old and crowded; we
have tried to record it accurately rather than minimally. The debt to
Levent Alp\"oge's example, and to Akhil Mathew for posing the problem, is
total. Computations were carried out with Claude (Fable~5), SymPy 1.14,
and msolve 0.10.1.


\subsection*{Note added (August 4, 2026)}
A final pre-release survey found two independent observations adjacent to Theorem~C:
a note of A.~Husain (github.com/Cobord/Jacobian, August 1, 2026) constructs the
cotangent-lift Poisson endomorphism and observes the intertwining of the weighted tori,
posing the isotropy computation as open, without moment-map structure; and a blog comment
of D.~Fillmore (July 31, 2026, on Tao's digestion post) interprets the $3$-to-$1$ fiber
via Majorana spin addition, again without symplectic content. Neither states the
moment-map identity $\nu\circ\Phi=\mu$ proved here; we record them for completeness.

\begin{thebibliography}{99}
\bibitem{Hus26} A.~Husain, \emph{Jacobian--Dixmier monoid} (working note), \texttt{github.com/Cobord/Jacobian}, August 1, 2026.
\bibitem{Fil26} D.~Fillmore, comment on T.~Tao's digestion post, July 31, 2026.


\bibitem[AvdE07]{AvdE07}
P.~K. Adjamagbo and A. van den Essen,
\emph{A proof of the equivalence of the Dixmier, Jacobian and Poisson
conjectures},
Acta Math. Vietnam. \textbf{32} (2007), no. 2--3, 205--214;
arXiv:math/0608009.

\bibitem[Alp26]{Alpoge26}
L. Alp\"oge,
\emph{A counterexample to the Jacobian conjecture in dimension three},
public announcement, July 19, 2026. See \cite{Tao26, Speyer26} for
expositions.

\bibitem[BCW82]{BCW82}
H. Bass, E.~H. Connell, and D. Wright,
\emph{The Jacobian conjecture: reduction of degree and formal expansion
of the inverse},
Bull. Amer. Math. Soc. (N.S.) \textbf{7} (1982), 287--330.

\bibitem[Bav05]{Bavula05}
V.~V. Bavula,
\emph{The Jacobian conjecture${}_{2n}$ implies the Dixmier
problem${}_{n}$},
arXiv:math/0512250 (2005).

\bibitem[BL20]{BL20}
V.~V. Bavula and V. Levandovskyy,
\emph{A remark on the Dixmier conjecture},
Canad. Math. Bull. \textbf{63} (2020), no. 1, 6--12.

\bibitem[BKK05]{BKK05}
A. Belov-Kanel and M. Kontsevich,
\emph{Automorphisms of the Weyl algebra},
Lett. Math. Phys. \textbf{74} (2005), 181--199;
arXiv:math/0512169.

\bibitem[BKK07]{BKK07}
A. Belov-Kanel and M. Kontsevich,
\emph{The Jacobian conjecture is stably equivalent to the Dixmier
conjecture},
Mosc. Math. J. \textbf{7} (2007), no. 2, 209--218;
arXiv:math/0512171.

\bibitem[Cam73]{Campbell73}
L.~A. Campbell,
\emph{A condition for a polynomial map to be invertible},
Math. Ann. \textbf{205} (1973), 243--248.

\bibitem[Dix68]{Dixmier68}
J. Dixmier,
\emph{Sur les alg\`ebres de Weyl},
Bull. Soc. Math. France \textbf{96} (1968), 209--242.

\bibitem[Gao26]{Gao26}
S. Gao,
\emph{Counterexamples to the Jacobian conjecture in dimensions greater
than two},
arXiv:2608.00222v1, July 31, 2026.

\bibitem[HK26]{HuqKuruvilla26}
I. Huq-Kuruvilla,
\emph{An explicit characteristic-$2$ counterexample to the separable
Jacobian conjecture},
arXiv:2607.20968v1, July 23, 2026.

\bibitem[Jel26]{Jelonek26}
Z. Jelonek,
\emph{On mappings with Jacobian one},
arXiv:2607.20597v1, July 22, 2026.

\bibitem[Dru83]{Druzkowski83}
L.~M. Dru\.zkowski,
\emph{An effective approach to Keller's Jacobian conjecture},
Math. Ann. \textbf{264} (1983), 303--313.

\bibitem[KBEY18]{KBEY18}
A. Kanel-Belov, A. Elishev, and J.-T. Yu,
\emph{Automorphisms of Weyl Algebra and a Conjecture of Kontsevich},
arXiv:1802.01225 (2018).

\bibitem[Kel39]{Keller39}
O.-H. Keller,
\emph{Ganze Cremona-Transformationen},
Monatsh. Math. Phys. \textbf{47} (1939), 299--306.

\bibitem[May26]{Mayner26}
W.~G.~P. Mayner,
\emph{The Dixmier conjecture fails for $A_{n}$, $n\ge3$} (informal note,
6 pp., \texttt{dixmier-note.tex}) together with \emph{Structural analysis
of the Jacobian conjecture counterexample in $\CC^{3}$}
(\texttt{REPORT.md}), both prepared with Claude Fable~5,
\url{https://github.com/wmayner/dixmier-counterexample}, commits of
July 21, 2026, 05:36 and 06:41 UTC; the Dixmier statement was first made
in a comment on \cite{Speyer26}, July 20, 2026. Section references of the
form \S$n$ are to \texttt{REPORT.md}. The symplectic cotangent lift is
\S8 of \texttt{REPORT.md} (verified-facts table item~24; priority
list \S12, item~3) and item~2 of ``What this does not give'' in
\texttt{dixmier-note.tex}.

\bibitem[Lon26]{Long26}
C.~D. Long,
\emph{Small counterexamples to the Gaussian Moments conjecture},
arXiv:2607.18186v1, July 20, 2026.

\bibitem[Lou26]{Lou26}
A. Lou,
\emph{A derivation of the Jacobian conjecture counterexample},
note, \url{https://aaronlou.com/jacobian_counterexample_derivation.pdf};
PDF timestamp July 20, 2026.

\bibitem[MO26]{MO513387}
\emph{Galois structure of the new counterexample to the Jacobian
conjecture: an explicit cubic model with $S_{3}$ monodromy --- is this
known?},
MathOverflow question 513387, posted July 20, 2026, 07:56 UTC
(no answers as of August 4, 2026),
\url{https://mathoverflow.net/q/513387}.

\bibitem[MY26]{MengYang26}
G. Meng and L. Yang,
\emph{A five-variable counterexample to the Hessian conjecture, and the
low-dimensional status of the Jacobian and Hessian conjectures},
arXiv:2607.22198v2, July 24/27, 2026.

\bibitem[Mig26]{Migus26}
P. Migus,
\emph{Generic degrees of real polynomial Keller maps with non-dense
image},
arXiv:2607.21572v2, July 23/30, 2026.

\bibitem[Moh83]{Moh83}
T.~T. Moh,
\emph{On the Jacobian conjecture and the configurations of roots},
J. Reine Angew. Math. \textbf{340} (1983), 140--212.

\bibitem[Sha26]{Shaska26}
T. Shaska,
\emph{Graded Keller maps and the Jacobian Conjecture},
arXiv:2607.20210; v1 July 22, 2026, v2 July 25, 2026. \emph{References
here are to v2 unless stated otherwise; Thm.~7.1 of v2 is Thm.~6.1 of v1,
and Thm.~8.3, Lem.~8.5, \S9.2 and \S10 appear only in v2.}

\bibitem[Spe26]{Speyer26}
D.~E. Speyer,
\emph{The new counterexample to the Jacobian conjecture},
blog post with comment thread, Secret Blogging Seminar, July 20, 2026,
\url{https://sbseminar.wordpress.com/2026/07/20/the-new-counterexample-to-the-jacobian-conjecture/}.

\bibitem[Tao26]{Tao26}
T. Tao,
\emph{A digestion of the Jacobian conjecture counterexample},
blog post, What's new, July 21, 2026,
\url{https://terrytao.wordpress.com/2026/07/21/a-digestion-of-the-jacobian-conjecture-counterexample/}.

\bibitem[Tsu05]{Tsuchimoto05}
Y. Tsuchimoto,
\emph{Endomorphisms of Weyl algebra and $p$-curvatures},
Osaka J. Math. \textbf{42} (2005), no. 2, 435--452.

\bibitem[Ula26]{Ulam26}
Anonymous,
\emph{A counterexample to the Jacobian conjecture},
note (7 pp.), \url{https://ulam.ai/research/jacobian.pdf}; PDF timestamp
July 20, 2026, 05:54 EDT.

\bibitem[vdE00]{vdE00}
A. van den Essen,
\emph{Polynomial Automorphisms and the Jacobian Conjecture},
Progress in Mathematics \textbf{190}, Birkh\"auser, Basel, 2000.

\bibitem[Wan80]{Wang80}
S.~S.-S. Wang,
\emph{A Jacobian criterion for separability},
J. Algebra \textbf{65} (1980), no.~2, 453--494.

\bibitem[Yag80]{Yagzhev80}
A.~V. Yagzhev,
\emph{On Keller's problem},
Sibirsk. Mat. Zh. \textbf{21} (1980), no.~5, 141--150; English transl.,
Siberian Math. J. \textbf{21} (1980), 747--754.

\bibitem[Zhu26]{Zhu26}
B. Zhu,
\emph{Non-injective field redefinitions and quantum inequivalence in
scalar theories},
arXiv:2607.18166v1, July 20, 2026.

\bibitem[Zhe24]{Zheglov24}
A.~B. Zheglov,
\emph{The Conjecture of Dixmier for the first Weyl algebra is true},
arXiv:2410.06959 (2024; v5, January 19, 2026). Claimed proof of
$\DC_{1}$; under review at the time of writing.

\end{thebibliography}

\end{document}
